\section{Classical inputs and Theorem 2.8}\label{p1:sec:classical-inputs}

Put $e(t)=\exp(2\pi i t)$. Write $n\sim N$ for $N<n\leq2N$. Unspecified arithmetic indices are positive integers. Norms of sequences are counting-measure $\ell^2$ norms. A star denotes reduced residue classes in residue sums. In character sums it denotes primitive characters. These two uses are distinguished explicitly below. Constants are uniform within each sequence family. Put
\[
 \theta_{\mathrm{KS}}=\frac7{32}.
\tag{2.1}\label{p1:eq:2.1}
\]
Thus $\theta_{\mathrm{KS}}=2\cdot7/64$, where $7/64$ is the unramified Ramanujan exponent. This symbol is distinct from the level parameter $\theta$ and the exceptional parameters $\vartheta_j$.

\subsection{Spectral normalization and large sieves}

The spectral large sieves are due to Deshouillers and Iwaniec \cite{DI82}.

Let $q\geq1$. Use Petersson-orthonormal bases at level $\Gamma_0(q)$. The measure is $dx\,dy/y^2$. Take $\sigma_\infty=I$. For a Maass cusp form,
\[
 u_j(z)=\sum_{n\ne0}\rho_{j\infty}(n)\sqrt y\,
 K_{i\kappa_j}(2\pi|n|y)e(nx),\qquad
 \lambda_j=\tfrac14+\kappa_j^2.
\tag{2.2}\label{p1:eq:2.2}
\]
This is the Fourier expansion of \cite[(1.15)]{DI82}, as printed there.
It has no factor $2\sqrt{|n|}$; that factor enters only through the dictionary (2.11).
For a holomorphic form $f$ of even weight $k\geq2$, write
\[
 f(z)=\sum_{n\geq1}\psi_f(\infty,n)e(nz).
\tag{2.3}\label{p1:eq:2.3}
\]
The Eisenstein coefficients $\varphi_{\mathfrak c\infty n}(s)$ have the normalization of \cite[(1.17)]{DI82}. Their relation to the Whittaker coefficients $\rho^{\rm Dr}$ of Drappeau, whose full dictionary is (2.11) in \S2.2 below, is
\[
 \rho^{\rm Dr}_{\mathfrak c\infty}(n,t)
 =\frac{\pi^{1/2+it}|n|^{it-1/2}}{\Gamma(1/2+it)}
   \varphi_{\mathfrak c\infty n}(1/2+it).
\tag{2.4}\label{p1:eq:2.4}
\]
Let $\mathfrak c$ run over inequivalent cusps. Use width-one scaling matrices. The continuous-spectrum measure in the DI large sieve is plain $dt$.

\textbf{Theorem 2.1 (consequence of \cite[Theorem~2]{DI82}).} Let $\epsilon>0$, $q\geq1$, $K\geq1$ and $N\geq\tfrac12$. Let $a_n\in\mathbb C$ be supported on $N<n\leq2N$. Then each of
\[
\begin{split}
 H(K,a)&=\sum_{\substack{2\leq k\leq K\\k\in{\rm even}}}
 \frac{(k-1)!}{(4\pi)^{k-1}}
 \sum_{f\in\mathcal B_k(q)}
 \left|\sum_{n\sim N}a_n n^{-(k-1)/2}\psi_f(\infty,n)\right|^2,\\
 M(K,a)&=\sum_{|\kappa_j|\leq K}
 \frac1{\cosh(\pi\kappa_j)}
 \left|\sum_{n\sim N}a_n\rho_{j\infty}(n)\right|^2,\\
 E(K,a)&=\sum_{\mathfrak c}\int_{-K}^{K}
 \left|\sum_{n\sim N}a_n n^{it}
 \varphi_{\mathfrak c\infty n}(1/2+it)\right|^2dt
\end{split}
\tag{2.5}\label{p1:eq:2.5}
\]
is at most
\[
 C_\epsilon\left(K^2+\frac{N^{1+\epsilon}}q\right)
 \sum_{n\sim N}|a_n|^2,
\tag{2.6}\label{p1:eq:2.6}
\]
uniformly in $q,K,N,a$. By \cite[Theorem~2, (1.28)--(1.30)]{DI82},
each expression is bounded by
$C_\epsilon\{K^2+\mu(\mathfrak a)N^{1+\epsilon}\}\|a_N\|_2^2$,
with the constant depending only on $\epsilon$. At
$\mathfrak a=\infty$, $\mu(\infty)=1/q$, giving (2.6); the source includes
the endpoint $N\geq\tfrac12$. The third line of (2.5) has no factor
$1/\cosh(\pi t)$. The cusp transport for Theorem 2.8 is in
\S\ref{p1:sec:t-cusp-transport}.

For an exceptional eigenvalue $0<\lambda_j<\tfrac14$, choose $i\kappa_j>0$, and put
\[
 \vartheta_j=2i\kappa_j=2\sqrt{\tfrac14-\lambda_j}.
\tag{2.7}\label{p1:eq:2.7}
\]

\textbf{Theorem 2.2 (\cite[Theorem~5]{DI82}, the specialization used here).} Let $\epsilon>0$, $q\geq1$, $N\geq\tfrac12$ and $Y\geq1$. Let $a_n\in\mathbb C$ be supported on $n\sim N$. Then
\[
 \sum_{j\ {\rm exceptional}}Y^{\vartheta_j}
 \left|\sum_{n\sim N}a_n\rho_{j\infty}(n)\right|^2
 \ll_\epsilon
 \left(1+\sqrt{\frac{NY}q}\right)
 \left(1+\sqrt{\frac{N^{1+\epsilon}}q}\right)\|a\|_2^2.
\tag{2.8}\label{p1:eq:2.8}
\]
This is \cite[Theorem~5]{DI82}. It uses the notation of Theorem 2, so $N\geq\tfrac12$ is allowed. The implied constant depends on $\epsilon$ alone.

\textbf{Theorem 2.3 (from \cite[proof of Theorem~7, (8.17)]{DI82}).} Let $\epsilon>0$. Let $a_n$ be the characteristic sequence of an integer interval
$I\subset(N,2N]$. Then, for some $c_\epsilon$ and uniformly in $Q,N,Y\ge1$ and $I$,
\[
\sum_{2\le q\le Q}\ \sum_{j\ {\rm exceptional}}^{(q)}Y^{\vartheta_j}
\Bigl|\sum_{n\sim N}a_n\rho_{j\infty}(n)\Bigr|^2\le c_\epsilon(QN)^\epsilon\bigl(Q+N+\sqrt{NY}\bigr)N.
\tag{2.9}\label{p1:eq:2.9}
\]
The range starts at $q=2$ because Theorem 2.8 uses level $q=rs$ with $r>R\ge1$, so $q\ge2$ there.

\emph{Proof.} By \cite[(8.17) and (8.4)]{DI82},
\[
\begin{aligned}
&\sum_{Q_0<q\le16Q_0}\sum_{\rm exc}Y^{\vartheta_j}
 \left|\sum_{N<n\le2N}a_n\rho_{j\infty}(n)\right|^2\\
&\qquad\le c^{\rm DI}_{\epsilon'}(Q_0N)^{5\epsilon'}
                  (Q_0+N+\sqrt{NY})N
\end{aligned}
\]
for $Q_0\ge1$. In \cite{DI82} the interval is printed as $[N,N_1]$, $N\le N_1\le2N$. The sum runs over $N<n\le2N$, so $a_n$ is the characteristic sequence of the initial interval
$(N,N_1]\cap\mathbb Z$. The source estimate (8.17) is uniform for
$Q_0,N,Y\ge1$ and every real $N_1\in[N,2N]$; its final statement imposes
no relation between $Y$ and $Q_0,N$. In the proof, the induction first treats
$1\le N\le Q_0$ at $Y_*=Q_0^{2-2\epsilon'}/N$, while (8.18) handles the
remaining large-$N$ range, and the final comparison extends the result to
every $Y\ge1$. We write all intervals in this half-open form.

(i) \emph{$q$-blocks.} Let $I_Q=\max\{i\ge0:16^{-i}Q\ge1\}$. The blocks $(16^{-i}Q,16^{1-i}Q]$,
$1\le i\le I_Q$, together with $(1,16]$, cover $2\le q\le Q$. There are at most $2+\log_{16}Q$ of
them. Every summand is nonnegative. So each block may be enlarged to its full DI range. Each block has
$Q_0\le Q$.

(ii) \emph{Intervals.} Write $I=(u,v]\cap\mathbb Z$ with $N\le u\le v\le2N$. Then
$\mathbf1_I=\mathbf1_{(N,v]}-\mathbf1_{(N,u]}$, both initial intervals in DI's sense, and
$|A-B|^2\le2|A|^2+2|B|^2$. This costs a factor 4.

Hence the left side of (2.9) is at most
$4c^{\rm DI}_{\epsilon'}(2+\log_{16}Q)(QN)^{5\epsilon'}(Q+N+\sqrt{NY})N$. Take $\epsilon'=\epsilon/10$
and $c_\epsilon:=4c^{\rm DI}_{\epsilon/10}\sup_{Q\ge1}(2+\log_{16}Q)Q^{-\epsilon/2}<\infty$. \hfill$\square$

The interval hypothesis is the one in \cite[(8.17)]{DI82}, and the induction closes in the proof of Theorem~7. Theorem 2.3 comes from this step of that proof. \cite[Theorem~7]{DI82} treats $a_n=1$ on $n\le N$. Theorem 2.3 needs $N\ge1$. The one-integer block $N=\tfrac12$ is
treated in \S\ref{p1:sec:t-exceptional}.

\subsection{Hecke exponents and the adopted trace normalization}

\textbf{Kim--Sarnak input.} The bound $7/64$ at an unramified place is \cite[Appendix~2, Proposition~2]{KS03}; the Satake normalization is stated in \cite[p.~175]{KS03}. Independently, the archimedean input
$\lambda_j\geq975/4096$ is the lower bound in \cite[Appendix~2]{KS03}; it is not a consequence of the finite-place estimate below. Let $f$ be a weight-zero Maass newform of congruence level. Each exceptional parameter of $f$ satisfies $0\leq i\kappa\leq7/64$. For a prime $p$ not dividing the level,
% Locator retained: the Satake normalization passage on printed p. 175 has no theorem or equation number.
\[
 |\alpha_p|,|\beta_p|\leq p^{7/64},\qquad
 |\lambda_f(p^\nu)|\leq(\nu+1)p^{7\nu/64}
 \quad(\nu\geq0).
\tag{2.10}\label{p1:eq:2.10}
\]
Thus the preceding archimedean input gives $\vartheta_j\leq\theta_{\mathrm{KS}}$, equivalently $\lambda_j\geq975/4096$ by (2.7). At the finite places, (2.10) gives $|\lambda_f(k)|\leq\tau(k)k^{\theta_{\mathrm{KS}}/2}$ for $k$ coprime to the level. We use these bounds for newforms of level dividing $q$. These newforms diagonalize the unramified Hecke operators on the full level-$q$ space. For holomorphic newforms we use Deligne's $ |\lambda_f(k)|\leq\tau(k)$, again for $k$ coprime to the level. Both factorizations use primes coprime to the level.

\textbf{Deligne input.} The holomorphic bound is \cite[Th\'eor\`eme~8.2]{Deligne74}. Both inputs enter at primes coprime to the level, through the prime-power recurrences and multiplicativity above.

Bruggeman and Kuznetsov independently introduced the relevant trace
formulas; see \cite[p.~220]{DI82} and their original papers
\cite{Bruggeman78,Kuznetsov81}. For the normalization and formulas used
here, see \cite[(4.7), (4.9)--(4.19), Lemma~4.5]{Drappeau}, in the PLMS version.
% Locator retained: DI82's historical attribution on printed p. 220 has no theorem or equation number.
The coefficient dictionary is
\[
\begin{aligned}
 \rho^{\rm Dr}_{f\mathfrak a}(n)&=\frac{\rho^{\rm DI}_{f\mathfrak a}(n)}{2\sqrt{|n|}}
 &&\hbox{(Maass)},\\
 \rho^{\rm Dr}_{f\mathfrak a}(n)&=(4\pi n)^{-k/2}\psi_f(\mathfrak a,n)
 &&\hbox{(holomorphic)},
\end{aligned}
\tag{2.11}\label{p1:eq:2.11}
\]
with (2.4) for Eisenstein series. The Maass identity follows from $W_{0,it}(4\pi|n|y)=2\sqrt{|n|y}K_{it}(2\pi|n|y)$.

For $\phi\in C_c^\infty((0,\infty))$, define the DI transforms
\[
\begin{split}
 \widehat\phi(t)&=\frac\pi{\sinh\pi t}
 \int_0^\infty\frac{J_{2it}(x)-J_{-2it}(x)}{2i}\phi(x)\frac{dx}x,\\
 \check\phi(t)&=\frac4\pi\cosh\pi t
 \int_0^\infty K_{2it}(x)\phi(x)\frac{dx}x,\\
 \widetilde\phi(k-1)&=\int_0^\infty J_{k-1}(x)\phi(x)\frac{dx}x.
\end{split}
\tag{2.12}\label{p1:eq:2.12}
\]
Values at removable singularities are limits. In Drappeau's notation the corresponding transforms are $-4\widehat\phi$, $2\pi\check\phi$, and $4i^k\widetilde\phi(k-1)$.

Theorem 2.8 and Section~3 use the following identity. Let $q=rs$, $(r,s)=1$, and $r>1$; take
\[
 \sigma_{1/s}=\begin{pmatrix}\sqrt r&0\\s\sqrt r&1/\sqrt r\end{pmatrix},
 \qquad
 \mathcal C(\infty,1/s)=\{\gamma=cs\sqrt r:c\geq1, (c,r)=1\}.
\tag{2.13}\label{p1:eq:2.13}
\]
The scaling matrix and admissible moduli in (2.13) are also given in
\cite[(1.6) and (2.4)]{DI82} and \cite[Lemma~4.3]{Drappeau}.
Since $r>1$, the cusps $\infty$ and $1/s$ are inequivalent. Hence the same-sign identity has no diagonal term. If a matrix in $\Gamma_0(rs)$ maps $\infty$ to the reduced fraction $1/s$, its lower-left entry is $\pm s$. Then $rs\mid s$, so $r=1$. Every application here has $r>R\geq1$. For $m,n\geq1$, put $\mathcal K_\pm(m,n)=\sum_{\gamma\in\mathcal C}\gamma^{-1}S_{\infty,1/s}(m,\pm n;\gamma)\phi(4\pi\sqrt{mn}/\gamma)$. In the DI coefficients, Drappeau's formulas read
\[
\begin{split}
 \mathcal K_+(m,n)
 &=\frac1\pi\sum_{\substack{k\geq2\\k\ {\rm even}}}
 i^k(k-1)!(4\pi\sqrt{mn})^{1-k}\widetilde\phi(k-1)
 \sum_f\overline{\psi_f(\infty,m)}\psi_f(1/s,n)\\
 &\quad-\sum_j\frac{\widehat\phi(\kappa_j)}{\cosh\pi\kappa_j}
 \overline{\rho_{j\infty}(m)}\rho_{j,1/s}(n)\\
 &\quad-\frac1\pi\sum_{\mathfrak c}\int_{\mathbb R}
 \widehat\phi(t)(n/m)^{it}
 \overline{\varphi_{\mathfrak c\infty m}(1/2+it)}
 \varphi_{\mathfrak c,1/s,n}(1/2+it)\,dt,\\[2mm]
 \mathcal K_-(m,n)
 &=\frac\pi2\sum_j\frac{\check\phi(\kappa_j)}{\cosh\pi\kappa_j}
 \overline{\rho_{j\infty}(m)}\rho_{j,1/s}(-n)\\
 &\quad+\frac12\sum_{\mathfrak c}\int_{\mathbb R}
 \check\phi(t)(n/m)^{it}
 \overline{\varphi_{\mathfrak c\infty m}(1/2+it)}
 \varphi_{\mathfrak c,1/s,-n}(1/2+it)\,dt.
\end{split}
\tag{2.14}\label{p1:eq:2.14}
\]
The proof uses only the identities (2.7.T16)--(2.7.T18), derived in \S\ref{p1:sec:t-cusp-pair}. The continuous integrals carry plain $dt$. Let $\alpha_m,\beta_n$ be finite sequences. Multiply (2.14) by $\alpha_m\beta_n$, sum in $(m,n)$, then take spectral absolute values and apply Cauchy--Schwarz. The first functional is $L_1(u)=\sum_m\overline{\alpha_m}\rho_{u\infty}(m)$. The spectral product is $\overline{L_1(u)}\sum_n\beta_n\rho_{u,1/s}(\pm n)$.

\subsection{BFI and the classical progression estimates}

For a finitely supported sequence $f$, define
\[
 \Delta_f(q,a)=\sum_{n\equiv a\ (q)}f(n)
 -\frac1{\varphi(q)}\sum_{(n,q)=1}f(n),\qquad (a,q)=1.
\tag{2.15}\label{p1:eq:2.15}
\]
The notation of \cite[(1.1)]{BFI86} writes $\Delta_f(x;q,a)$ for (2.15) with both sums cut at $n\le x$.

The progression estimates are from Bombieri, Friedlander and Iwaniec
\cite{BFI86}. We use \cite[(A2), Theorem~0(a), and proof of Theorem~0(a)]{BFI86}.

\textbf{BFI hypothesis (A2) \cite[(A2)]{BFI86}.} Fix $B_\beta>0$ for the sequence family. Let $\beta$ be in the family, with support parameter $N\geq2$. For every $d,k\geq1$ and integer $\ell\ne0$ with $(\ell,k)=1$, and for every $A>0$, there is $C_A$ with
\[
 \left|\sum_{\substack{n\equiv\ell\ (k)\\(n,d)=1}}\beta_n
 -\frac1{\varphi(k)}\sum_{(n,dk)=1}\beta_n\right|
 \leq C_A\|\beta\|_2N^{1/2}\tau(d)^{B_\beta}(\log(2N))^{-A}.
\tag{2.16}\label{p1:eq:2.16}
\]
The $C_A$ are independent of $N,d,k,\ell$ and of $\beta$ in the family. As in \cite{BFI86}, $d$ and $k$ run over all positive integers. If $\|\beta\|_2^2\geq N(\log 2N)^{-B'}$ for a fixed $B'$, then a bound $C_AN\tau(d)^{B_\beta}(\log 2N)^{-(A+B'/2)}$ implies (2.16), after adjusting the constant.

\textbf{Theorem 2.4 (\cite[Theorem~0(a), (2.2)]{BFI86}).} Let $\beta_n$ be supported on $n\leq N$, $N\geq2$, and satisfy (2.16). For every $A>0$, there is $B_1>0$ such that, uniformly for $Q\leq N(\log N)^{-B_1}$,
\[
 \sum_{q\leq Q}\ \sum_{\substack{a\bmod q\\(a,q)=1}}
 |\Delta_\beta(q,a)|^2
 \ll \|\beta\|_2^2N(\log N)^{-A}.
\tag{2.17}\label{p1:eq:2.17}
\]
The constants depend on $A,B_\beta$ and the constants in (2.16). For a sequence on $N<n\leq2N$, replace $N$ by $2N$ in (2.17).

\cite[Theorem~0(b)]{BFI86} is not used; its discrepancy has the cutoff $n\leq x$ of \cite[(1.1)]{BFI86}, while (2.15) is a full-support discrepancy. The switched bilinear form uses the full-support bound (3.49), proved from (3.25), (3.48) and Theorem 2.4.

That bound uses \cite[(1.6)]{BFI86}. Let $c_h\in\mathbb C$ be supported on $1\leq h\leq H$. The large sieve gives
\[
 \sum_{q\leq Q}
 \sum_{\chi\bmod q}^{\rm primitive}
 \left|\sum_hc_h\chi(h)\right|^2
 \ll(Q^2+H)\|c\|_2^2.
\tag{2.19}\label{p1:eq:2.19}
\]
We use this initial-interval form with $H$ the upper support bound.

\textbf{Bombieri--Vinogradov input.} Write $\psi(y;q,a)=\sum_{n\leq y,n\equiv a\ (q)}\Lambda(n)$. For every $A>0$ there is $B=B(A)>0$ such that, for $y\geq3$ and $Q\leq y^{1/2}(\log y)^{-B}$,
\[
 \sum_{q\leq Q}\max_{(a,q)=1}
 \left|\psi(y;q,a)-\frac y{\varphi(q)}\right|
 \ll_A y(\log y)^{-A}.
\tag{2.20}\label{p1:eq:2.20}
\]
This is BFI's formulation of the Bombieri--Vinogradov theorem
\cite[(1.4)]{BFI86}; compare the original results
\cite{Bombieri65,Vinogradov65}. The maximum allows the varying reduced
residue $a=N\bmod q$. We apply (2.20) at each endpoint separately, with
its own level inequality. Thus (2.20) is used at fixed $y$ only.

\textbf{Siegel--Walfisz input.} We use the power-log form of
\cite[Corollary~5.29]{IK04}. Remove prime powers and sum by parts.
Let $C>0$ and $A>0$ be fixed. For $y\geq3$, $q\leq(\log y)^C$
and $(a,q)=1$,
\[
 \pi(y;q,a)=\frac{\operatorname{Li}(y)}{\varphi(q)}
 +O_{A,C}\bigl(y(\log y)^{-A}\bigr).
\tag{2.21}\label{p1:eq:2.21}
\]
The constant is uniform in reduced $a$ and is ineffective. Subtraction at two comparable endpoints keeps any fixed logarithmic saving. We use (2.21) with this power-log error. Three further forms need their own proofs. The first saves a power of the interval length. The second keeps the tensor and coprimality structure of the structured sequences. The third is localized to one dyadic box of the Heath--Brown decomposition. Section~3 proves the short-interval estimates and those for its structured sequences.

\subsection{The linear sieve and its structured remainder}

Let $\mathcal A$ be a finite integer sequence with nonnegative multiplicities. Let $\mathcal P$ be a prime set. Put
\[
 P(z)=\prod_{\substack{p<z\\p\in\mathcal P}}p,\quad
 S(\mathcal A,\mathcal P,z)=\sum_{a\in\mathcal A}1_{(a,P(z))=1},
 \quad |\mathcal A_d|=\frac{\omega(d)}dX+r(\mathcal A,d).
\tag{2.22}\label{p1:eq:2.22}
\]
Here $X>0$ and $\omega$ is multiplicative. Require $0<\omega(p)<p$ for $p\in\mathcal P$. Set $\omega(p)=0$ for excluded primes. Set $V(z)=\prod_{p<z,p\in\mathcal P}(1-\omega(p)/p)$. Use the continuous linear-sieve solutions
\[
\begin{gathered}
 sF(s)=2e^\gamma\quad(0<s\leq3),\qquad
 sf(s)=0\quad(0<s\leq2),\\
 (sF(s))'=f(s-1)\quad(s>3),\qquad
 (sf(s))'=F(s-1)\quad(s>2).
\end{gathered}
\tag{2.23}\label{p1:eq:2.23}
\]
On $s\geq1$, $F$ decreases to 1. Also $f$ increases to 1, with $f=0$ through $s=2$. In particular $F(1)=2e^\gamma$. We use this endpoint when the level is at least the sifting limit.

Let $K_{\mathrm M}\geq2$ be one fixed absolute constant such that, for all
$z>w\geq2$,
\[
 \prod_{w\leq p<z}\left(1-\frac1p\right)^{-1}
 \leq\frac{\log z}{\log w}\left(1+\frac{K_{\mathrm M}}{\log w}\right).
\]
Such a constant exists by Mertens' product estimate, after increasing it
to cover bounded $w$.

\textbf{Theorem 2.5 (\cite[Theorem~1]{AA36}, dimension one).} Suppose, for some fixed $K\geq2$ and every $z>w\geq2$,
\[
 \prod_{\substack{w\leq p<z\\p\in\mathcal P}}
 (1-\omega(p)/p)^{-1}
 \leq\frac{\log z}{\log w}\left(1+\frac K{\log w}\right).
\tag{2.24}\label{p1:eq:2.24}
\]
Let $y\geq z\geq2$ and $s=\log y/\log z$. Then
\[
\begin{split}
 S(\mathcal A,\mathcal P,z)
 &\leq XV(z)\{F(s)+e^{\sqrt K}Q_1(s)(\log y)^{-1/3}\}
 +\sum_{\substack{d<y\\d\mid P(z)}}|r(\mathcal A,d)|,\\
 S(\mathcal A,\mathcal P,z)
 &\geq XV(z)\{f(s)-e^{\sqrt K}Q_1(s)(\log y)^{-1/3}\}
 -\sum_{\substack{d<y\\d\mid P(z)}}|r(\mathcal A,d)|,
\end{split}
\tag{2.25}\label{p1:eq:2.25}
\]
where, for $s\geq1$,
\[
 Q_1(s)\leq\exp\{-s\log s+s\log\log(3s)+O(s)\}.
\tag{2.26}\label{p1:eq:2.26}
\]
See \cite{AA36}. The hypotheses are (1.1)--(1.3), and Theorem 1 is (1.4)--(1.6). The dimension-one constants $\beta=2$ and $A=2e^\gamma$ appear in \cite[p.~176]{AA36}; the source states $B=0$ in \cite[p.~173]{AA36}. For $s\leq2$ the lower main function in (2.25) vanishes. Hence its lower main term is nonpositive.
% Locators retained: the dimension-one table on p. 176 and the B=0 statement on p. 173 are unnumbered.

For the later prime set $\mathcal P_N=\{p:p\nmid2N\}$, take the
multiplicative function with $\omega(p)=1$ on $\mathcal P_N$,
$\omega(p)=0$ off it, and $\omega(p^a)=0$ for $a\geq2$. Its product
in (2.24) is over a subset of all primes in $[w,z)$, so (2.24) holds
uniformly in $N$ with the fixed value $K=K_{\mathrm M}$. The same data
satisfy (2.27) with $L_\omega=2$. Thus, whenever (2.25) is specialized
to this set, its factor is $e^{\sqrt{K_{\mathrm M}}}$. This is fixed;
combining it with (2.26) changes only the absolute constant in the
$O(s)$ term there.
This $\omega$ is a nominal sieve density: by Theorem~2.6 the constructed
coefficients are independent of $\mathcal A$, $X$, and $\omega$.

The structured remainder is from \cite[Theorem~4]{AA37}.

Assume (2.24) and nonnegativity of $\omega$ on prime powers. Fix $L_\omega>1$. Assume, for all $z>w\geq2$,
\[
 \sum_{\substack{w\leq p<z\\p\in\mathcal P}}
       \sum_{a\geq2}\frac{\omega(p^a)}{p^a}
 \leq\frac{L_\omega}{\log(3w)}.
\tag{2.27}\label{p1:eq:2.27}
\]
These are \cite[(1)--(2)]{AA37} in dimension one, with $L=L_\omega$.

Fix $0<\epsilon_s<1/3$ and $D\geq2$. Put $\delta_s=\epsilon_s^9$ and
\[
 \mathcal G(D)=\{D^{\epsilon_s^2(1+\delta_s)^n}:n\geq0\}.
\tag{2.28}\label{p1:eq:2.28}
\]
Let $\mathcal H(D)$ consist of tuples $\mathbf D=(D_1,\ldots,D_\rho)$ from this grid with $D_\rho\leq\cdots\leq D_1<D^{1/2}$. Define
\[
\begin{split}
 \mathcal D^+(D)&=\{\mathbf D\in\mathcal H(D):
 D_1\cdots D_{2l}D_{2l+1}^3<D
 \quad(0\leq l\leq(\rho-1)/2)\},\\
 \mathcal D^-(D)&=\{\mathbf D\in\mathcal H(D):
 D_1\cdots D_{2l-1}D_{2l}^3<D
 \quad(1\leq l\leq\rho/2)\}.
\end{split}
\tag{2.29}\label{p1:eq:2.29}
\]
\cite[Theorem~4]{AA37} gives $|\mathcal D^+|,|\mathcal D^-|<\exp(8\epsilon_s^{-3})$. This bound follows the definitions of $\mathcal D^+$ and $\mathcal D^-$. The grid is built from the relative level of each application.

\textbf{Theorem 2.6 (from \cite[Theorem~4]{AA37}, with $R=1$).} Assume (2.22), (2.24), (2.27), with the parameters above. Then there are coefficients $\phi_d^\pm$ and $\lambda_d^\pm(\mathbf D)$ of modulus at most 1. They are independent of $\mathcal A,X,\omega$. They satisfy
\[
\begin{aligned}
 S(\mathcal A,\mathcal P,z)&\leq XV(z)(F(s)+E_s)+R_1^++R^+
 &&(2\leq z\leq D),\\
 S(\mathcal A,\mathcal P,z)&\geq XV(z)(f(s)-E_s)-R_1^--R^-
 &&(2\leq z\leq D^{1/2}),
\end{aligned}
\tag{2.30}\label{p1:eq:2.30}
\]
where $s=\log D/\log z$ and
\[
 E_s\ll\epsilon_s+\epsilon_s^{-8}e^{K+L_\omega}(\log D)^{-1/3},
\tag{2.31}\label{p1:eq:2.31}
\]
with an absolute implied constant. \cite[Theorem~4]{AA37} contains $e^{K+L}$. For the coefficient construction used here, the prime set and multiplicative function just specified give $K=K_{\mathrm M}$ and $L_\omega=2$. Hence (2.31) has the fixed factor $e^{K_{\mathrm M}+2}$; it is retained in the formula and may be absorbed into the implied constant. The prime-power dependence comes from its (2), namely (2.27). In its estimate (31), put $u=D^{\epsilon_s^2}$. Then
\[
 U=\sum_{\substack{u\leq p<z\\p\in\mathcal P}}
       \sum_{a\geq2}\frac{\omega(p^a)}{p^a}\leq L_\omega,
 \qquad
 \prod_{\substack{u\leq p<z\\p\in\mathcal P}}
 \left\{\left(1-\frac{\omega(p)}p\right)
 \sum_{a\geq0}\frac{\omega(p^a)}{p^a}\right\}\leq e^U.
\]
For the first bound use $\max\{2,u\}$ in (2.27). An empty prime interval contributes zero. The second bound follows factor by factor. Thus $e^U\leq e^{L_\omega}$ bounds the extra prime-power factors in the nonsquarefree terms. Formula (2.31) keeps the factor $e^{K+L_\omega}$ of the source. The signed remainders are
\[
 R_1^\nu=\sum_{\substack{d<D^{\epsilon_s}\\d\mid P(D^{\epsilon_s^2})}}
 \phi_d^\nu r(\mathcal A,d),
\tag{2.32}\label{p1:eq:2.32}
\]
\[
 R^\nu=\sum_{\mathbf D\in\mathcal D^\nu(D)}
 \sum_{\substack{d<D^{\epsilon_s}\\d\mid P(D^{\epsilon_s^2})}}
 \lambda_d^\nu(\mathbf D)
 \sum_{\substack{D_i\leq p_i<D_i^{1+\delta_s}\\p_i\mid P(z)\\
                         p_i\ne p_j\ ((i,j)\in I(\mathbf D))}}
 r(\mathcal A,dp_1\cdots p_\rho).
\tag{2.33}\label{p1:eq:2.33}
\]
The sets $I(\mathbf D)$ impose distinctness on equal-label tail positions. Unequal grid intervals are disjoint. Coefficients may be zero on unused tuples. They depend only on $\mathcal P,D,z,\epsilon_s,\mathbf D$ and the distinctness conditions. In particular they are independent of $\mathcal A,X,\omega$. They may change with the level or the cutoff.

\cite[Theorem~4]{AA37} holds for general $R\geq1$, $0<\epsilon_s<3^{-R}$, with the respective ranges in (2.30). Its preceding combinatorial identities give $S\leq W_R^+$ and $S\geq W_R^-$. The coefficients come from \cite[(23)--(25), (26), and Lemma~4]{AA37}; these formulas also show their independence of the sequence.

\textbf{Relative-level grid.} Choose $0<\sigma\leq10^{-10}$ first, and put $\tau=\sigma^9$. For any level $L_c\ge2$, substitute
\[
 \epsilon_s=\sigma,\qquad D=L_c,\qquad
 \delta_s=\epsilon_s^9=\sigma^9=\tau.
\tag{2.33a}\label{p1:eq:2.33a}
\]
Then, term for term,
\[
 \mathcal G(D)
 =\{L_c^{\sigma^2(1+\tau)^n}:n\geq0\},\qquad
 D^{\epsilon_s^2}=L_c^{\sigma^2},\qquad
 D^{\epsilon_s}=L_c^\sigma.
\tag{2.33b}\label{p1:eq:2.33b}
\]
Thus the labels, half-open prime intervals, smooth-prime cutoff, and smooth-divisor support agree term by term with (2.32)--(2.33). These are the Iwaniec parameters at level $L_c$, and the tuple conditions (2.29) are unchanged. With this $\tau$, Theorem 2.6 applies as stated. In the application to (W), Section~5 takes $L_c=D=N^\theta$, so the grid there is exactly the global grid in Theorem~1.2. The coefficients are the values supplied by Iwaniec's construction for $(\mathcal P,L_c,\sigma,\mathbf D)$ at the cutoff $z$ of the application, here $z=N^u$. The cutoff also restricts the prime variables in (2.33). The proof uses only the bound $|\cdot|\le1$ and the support of these coefficients. It never compares coefficient values at different cutoffs.

\cite[Lemma~3]{AA37} gives $(F,f)$ with error
$O(e^{\sqrt K-s}(\log D)^{-1/3})$. In its use here the sieve
constant is $K=K_{\mathrm M}$, fixed independently of $N,D,z$.
Therefore $e^{\sqrt{K_{\mathrm M}}}$ is absorbed into the implied
constant, and the error is $O(e^{-s}(\log D)^{-1/3})$ with that fixed
constant. Likewise, in the profile (2.26), this multiplier changes only
its absolute $O(s)$ constant. The lemma cites
\cite[Lemmas~18 and 20]{AA36}. Theorem 4 includes the
closed endpoints $z=D$ and $z=D^{1/2}$ of (2.30), although
\cite[Lemma~20]{AA36} has strict inequalities.

\subsection{The fundamental lemma, Weil bound, and Heath--Brown identity}

\textbf{Fundamental lemma (cf. \cite[Lemma~4]{BFI86}; the
prime-support condition is explicit in
\cite[Lemma~5]{FI78}).} Let $D\geq2$, $s\geq3$ and
$z=D^{1/s}$. Let $P(z)$ be the product of all primes below $z$.
There are $\lambda_d^\pm$, supported on $d\leq D$, $d\mid P(z)$,
with $\lambda_1^\pm=1$, $|\lambda_d^\pm|\leq1$ and
\[
 (\lambda^-*1)(n)=(\lambda^+*1)(n)=1
 \quad\text{if }(n,P(z))=1,
\]
\[
 (\lambda^-*1)(n)\leq0\leq(\lambda^+*1)(n)
 \quad\text{otherwise},
\]
\[
\begin{aligned}
 \sum_{d\leq D}\frac{\lambda_d^\pm}d
 &=\prod_{p<z}(1-1/p)\\
 &\quad{}\times\{1+O(e^{-s})\}.
\end{aligned}
\tag{2.34}\label{p1:eq:2.34}
\]
The support $d\mid P(z)$ means that every supported $d$ is squarefree
and uses only primes below $z$. In (2.34), this makes
$(\lambda^\pm*1)(n)=\sum_{d\mid(n,P(z))}\lambda_d^\pm$ a sum over
exactly the sieving divisors: if $(n,P(z))=1$, only $d=1$ contributes,
and otherwise FI78's divisor-sum signs apply to $g=(n,P(z))$. The
support is also used in the remainder sum in (2.25). The explicit
Bonferroni weights in (3.102)--(3.104) use the same support to sieve
only primes in $\mathcal P$ below $z$ and to obtain the stated divisor
count. For the AA37 weights in (W), the corresponding condition
$d\mid P_N(Y)$ separates the smooth divisor, whose primes are below
$Y$, from tail primes $p_i\geq D_i\geq Y$. This is used in the unique
recovery of $d$ in the proof of the slot identity (4.13) and in
Appendix~\ref{app:proof-aa37}, Lemma~\ref{lem:aa37-support}.

The error $O(e^{-s})$ suffices below. Removing the primes of an auxiliary modulus changes the local Euler product and needs its own argument. The rough-integer estimates below use explicit Brun--Bonferroni weights; see also \cite{OperaCribro}. The Iwaniec endpoint $s=1$ follows from Theorem 2.6.

\textbf{Weil input.} For every $c\geq1$ and $u,v\in\mathbb Z$, define
\[
 S(u,v;c)=\sum_{x\bmod c,(x,c)=1}e((ux+v\bar x)/c).
\]
The required bound is
\[
 |S(u,v;c)|\leq\tau(c)(u,v,c)^{1/2}c^{1/2}.
\tag{2.35}\label{p1:eq:2.35}
\]
This is the Weil--Estermann bound with full gcd. Its gcd factor covers nonprimitive indices; see \cite[Corollary~11.12]{IK04}. In (2.14), (2.35) serves for absolute convergence. Each quantitative application keeps the gcd and modulus factors in its sums. For prime moduli this is Weil's bound; prime powers use Estermann's bound, and the general modulus follows by the Chinese remainder theorem.

\textbf{Lemma 2.7 (the exact Heath--Brown identity used here).} Let $X\geq2$, $K\geq1$ be an integer, and let $z\geq1$ satisfy $z^K>X$. Then for every integer $1\leq n\leq X$,
\[
 \Lambda(n)=\sum_{j=1}^{K}(-1)^{j-1}\binom Kj
 \sum_{\substack{m_1\cdots m_jn_1\cdots n_j=n\\m_i\leq z}}
 \mu(m_1)\cdots\mu(m_j)\log n_1.
\tag{2.36}\label{p1:eq:2.36}
\]
The identity is \cite[Lemma~1, (6)]{HB82}.

Equation (2.36) is algebraic. For $\Re s>1$, set $M_z(s)=\sum_{m\leq z}\mu(m)m^{-s}$, $Y(s)=\zeta(s)M_z(s)$. The coefficients of $Y-1$ vanish at indices $1\leq n\leq z$. The polynomial identity $1-(1-Y)^K=\sum_{j=1}^K(-1)^{j-1}\binom KjY^j$ gives
\[
 -\frac{\zeta'}\zeta
 =\sum_{j=1}^K(-1)^{j-1}\binom Kj
       (-\zeta')\zeta^{j-1}M_z^j
   +\left(-\frac{\zeta'}\zeta\right)(1-Y)^K.
\tag{2.37}\label{p1:eq:2.37}
\]
Each nonzero factor from $Y-1$ has index $>z$. Hence the last term has no coefficient of index $n\leq X$. Coefficient extraction proves (2.36). For the order-three decomposition on $n\leq2X_m$, take $K=3$, $z=2X_m^{1/3}$. Then $z^3=8X_m>2X_m$. The Möbius variables have support $m_i\leq2X_m^{1/3}$. Dyadic restrictions are imposed after this identity. Coefficient extraction preserves both product constraints and prime ordering.

\subsection{Theorem 2.8}

Theorem 2.8 belongs to the line of incomplete-Kloosterman estimates of
\cite[Theorem~12]{DI82} and
\cite[Theorem~2.1]{Drappeau}; Drappeau states that his theorem extends
DI82 Theorem~12. The closest refined quintilinear estimate is
\cite[Theorem~1.8, (1.8)--(1.9)]{Lichtman23}. Lichtman identifies
\cite{DPR23,ABL21} as antecedents for the exceptional-spectrum and
residue-dependence refinements. The statements are not identical: the
hypotheses (2.38)--(2.40) and bound (2.42) are the variant proved here
from (2.14) in \S\ref{p1:sec:theorem-2-8-proof}. Section~3 uses the full
bound (2.42) in (3.14)--(3.15).

\textbf{Theorem 2.8.}\label{p1:thm:T} Let $C,D,N,R,S\geq1$, $A\geq1$, and $0<\epsilon_0<1/100$. Let $b_{n,r,s}\in\mathbb C$ be supported on
\[
 0<n\leq N,\qquad R<r\leq2R,\qquad S<s\leq2S.
\tag{2.38}\label{p1:eq:2.38}
\]
Let $g:\mathbb R_{>0}^5\longrightarrow\mathbb C$ be smooth, supported in
\[
 (C,2C]\times(D,2D]\times\mathbb R_{>0}^3.
\]
For each multi-index $\nu=(\nu_1,\ldots,\nu_5)\in\mathbb Z_{\geq0}^5$, assume
\[
 |\partial^\nu g(c,d,n,r,s)|
 \leq G_\nu
 \bigl(c^{-\nu_1}d^{-\nu_2}n^{-\nu_3}r^{-\nu_4}s^{-\nu_5}\bigr)^{1-\epsilon_0}.
\tag{2.39}\label{p1:eq:2.39}
\]
The $G_\nu$ are uniform in the family. Let $a\in\mathbb Z$ satisfy $0<|a|\leq A$, and require
\[
 b_{n,r,s}=0\quad\text{unless }(a,rs)=1.
\tag{2.40}\label{p1:eq:2.40}
\]
Set $\mathcal X=2+ACDNRS$. For every $\epsilon>0$,
\[
\begin{split}
 &\left|\sum_{(r,s)=1}\sum_{0<n\leq N}b_{n,r,s}
 \sum_{\substack{c,d\geq1\\(rd,sc)=1}}
 g(c,d,n,r,s)e\!\left(\frac{an\overline{rd}}{sc}\right)\right|\\
 &\qquad\ll_{\epsilon,\epsilon_0,\{G_\nu\}}
 \mathcal X^{\epsilon+18\epsilon_0}J_*(a),
\end{split}
\tag{2.41}\label{p1:eq:2.41}
\]
where the inverse is modulo $sc$, and
\[
\begin{aligned}
 J_*(a)^2={}&
 \sum_{\substack{n''\mid a^\infty\\n''\leq2N}}
 A^{\theta_{\mathrm{KS}}}(n'')^{\theta_{\mathrm{KS}}}\\
 &\quad\times\Bigl[CS\Bigl(\frac N{n''}+RS\Bigr)(C+DR)
 +|a|NRS\Bigr]\|\widetilde b(n'')\|_2^2\\
 &+\Bigl[CS(CD\sqrt R)^{2\theta_{\mathrm{KS}}}
 (N+RS)^{1-\theta_{\mathrm{KS}}}\\
 &\qquad\times(C+DR)^{1-2\theta_{\mathrm{KS}}}
 +D^2NR\Bigr]\|b\|_2^2.
\end{aligned}
\tag{2.42}\label{p1:eq:2.42}
\]
Here $n''\mid a^\infty$ if each prime factor of $n''$ divides $a$. This includes $n''=1$. Also
\[
 \|b\|_2^2=\sum_{n,r,s}|b_{n,r,s}|^2,\qquad
 \|\widetilde b(n'')\|_2^2
 =\sum_{\substack{n',r,s\\(n',a)=1}}|b_{n'n'',r,s}|^2,
\tag{2.43}\label{p1:eq:2.43}
\]
with $b$ extended by zero outside (2.38). When $a=\pm1$, only $n''=1$ occurs in (2.42). The implied constant depends only on the displayed parameters and the seminorm bounds $G_\nu$. It is uniform in $(a,A,C,D,N,R,S,b)$.

For $(r,s)=1$, the condition $(rd,sc)=1$ means $(c,r)=1$ and $(d,sc)=1$. Pairs with $(c,s)>1$ are included, as are both signs of $a$. The sequences and weights may be complex.

\subsection{Proof of Theorem 2.8}\label{p1:sec:theorem-2-8-proof}

The proof uses the large-sieve bounds above, the Hecke bounds of Kim--Sarnak and Deligne, and the trace formula (2.14). We derive the opposite-sign trace formula below.

In this subsection $C,D,N_T,R_T,S_T,A$ are the parameters of Theorem 2.8. Write $R,S,N_1$ for $R_T,S_T,N_T$; Section~3 uses $R,S$ for other scales. Put
\[
\begin{aligned}
&\mathcal X=2+ACDN_1RS,\quad A_a=|a_T|,\quad q=rs,\\
&\quad(r,s)=1,
\quad\theta_{\mathrm{KS}}=7/32,\quad\vartheta_j=2i\kappa_j\in(0,\theta_{\mathrm{KS}}]
\end{aligned}
\tag{2.7.T1}\label{p1:eq:2.7.T1}
\]
on an exceptional eigenspace. All norms below use counting measure. Spectral bases are Petersson-orthonormal.

\subsubsection{Completion and separation}

Fix $n,r,s,c$, put $k=sc$, and define
$\widehat g_c(\xi)=\int g(c,t,n,r,s)e(-\xi t)\,dt$.
Group $d$ in reduced classes modulo $k$ and apply Poisson summation.
This gives the exact identity
\[
\sum_{(d,k)=1}g(c,d,n,r,s)e(a_Tn\overline{rd}/k)
 =k^{-1}\sum_{h\in\mathbb Z}\widehat g_c(h/k)
                       S(h\bar r,a_Tn;k).
\tag{2.7.T2}\label{p1:eq:2.7.T2}
\]
The zero frequency is bounded using
$|c_k(a_Tn)|\le(a_Tn,s)(a_Tn,c)$. Since
$\sum_{c\le2C}(u,c)\le2C\tau(u)$ and
$\sum_{s\le2S}(u,s)^2\le4S^2\tau(u)$, Cauchy--Schwarz gives its
absolute contribution
\[
\ll_\epsilon\mathcal X^\epsilon D(N_1R)^{1/2}\|b\|_2.
\tag{2.7.T3}\label{p1:eq:2.7.T3}
\]
This is the fourth term of (2.42).

Choose a small auxiliary exponent $\epsilon_1>0$. Set
\[
H=4SC D^{\epsilon_0-1}\mathcal X^{\epsilon_1},
\qquad\epsilon_2=\epsilon_1+\epsilon_0.
\tag{2.7.T4}\label{p1:eq:2.7.T4}
\]
Integrate by parts in $d$. Use (2.39) in $d$ on $d\in(D,2D]$. For every $\xi\ne0$ and $j\ge0$, this gives
$|\widehat g_c(\xi)|\ll_jD(|\xi|D^{1-\epsilon_0})^{-j}$. Every $h\ne0$ has $|h|/k\ge|h|/(4SC)$,
because $k=sc\le4SC$. By (2.7.T4) this gives $(|h|/k)D^{1-\epsilon_0}\ge(|h|/H)\mathcal X^{\epsilon_1}$.
Hence, for every $h$ with $|h|>H$,
\[
|\widehat g_c(h/k)|\ll_jD\,\mathcal X^{-j\epsilon_1}(H/|h|)^j .
\tag{2.7.T5}\label{p1:eq:2.7.T5}
\]
For $j\ge2$ and every $H>0$,
$\sum_{|h|>H}(H/|h|)^j\le2+2H$. If $H<1$, the sum runs over all $h\ne0$ and is at most
$2\zeta(2)H^2\le2+2H$. If $H\ge1$ it is at most $2(1+H/(j-1))$. With $|S(h\bar r,a_Tn;k)|\le k$,
the tail $|h|>H$ of (2.7.T2) is $\ll_j(1+H)D\mathcal X^{-j\epsilon_1}$ for each $(c,n,r,s)$. Summing
against $|b_{n,r,s}|$ over $c\le2C$ gives $\ll_j\mathcal X^{7-j\epsilon_1}\|b\|_2$, using
$\sum|b|\le(N_1RS)^{1/2}\|b\|_2$, $C,D,N_1,R,S\le\mathcal X$ and $1+H\le5\mathcal X^{1+\epsilon_1}$. Taking $j>20/\epsilon_1$ makes it smaller than (2.7.T3). If $H<1$, the retained
range $1\le|h|\le\lfloor H\rfloor$ is empty. Every nonzero frequency lies in this tail, including $h=\pm1$. Thus (2.7.T2) contains only the zero frequency and this tail. Equation (2.7.T3) completes this case.

For $1\le|h|\le\lfloor H\rfloor$, write
$h=\varsigma_hm$ and
$\varsigma=\varsigma_h\operatorname{sgn}(a_T)$, with signs
$\varsigma_h,\varsigma\in\{\pm1\}$.
The identity $S(-u,v;k)=S(u,-v;k)$ makes the second Kloosterman
argument $\varsigma A_an$. Partition the integers $m$ into
\[
\mathcal I_j=(M_j,\min(2M_j,\lfloor H\rfloor)]\cap\mathbb Z,
\qquad M_j=2^j/2.
\tag{2.7.T6}\label{p1:eq:2.7.T6}
\]
Thus the first nonempty block is exactly $\{1\}$, with $M=1/2$.
Partition $n$ into $(N',2N']$, $N'\ge1/2$. Shortened final
blocks remain integer intervals. There are $O(\log^2\mathcal X)$
pairs of blocks; each restricted coefficient norm is bounded by the
original norm.
For the current block define
\[
 b^\flat(n'')_{n',r,s}
 =b_{n'n'',r,s}\,1_{\{N'<n'n''\leq2N'\}}\,
  1_{\{(n',a_T)=1\}}.
\]
The slice satisfies $\|b^\flat(n'')\|_2\leq\|\widetilde b(n'')\|_2$.
Restrict $b$ to the current block. By support, $b^\flat(n'')=\widetilde b(n'')$ term by term. Thus the block estimate is exactly (2.42)--(2.43).

For one block write
\[
\begin{aligned}
\frac{\widehat g_c(\varsigma_hm/(sc))}{sc}
 &=\frac D{SC}W_0(c,m,n,r,s),\\
\gamma&=cs\sqrt r,\qquad W_1=W_0\gamma/(CS\sqrt R).
\end{aligned}
\tag{2.7.T7}\label{p1:eq:2.7.T7}
\]
Differentiate under the integral and use (2.7.T4). Then
\[
|\partial^\nu W_1|
 \ll_\nu\mathcal X^{|\nu|\epsilon_2}
            c^{-\nu_1}m^{-\nu_2}n^{-\nu_3}r^{-\nu_4}s^{-\nu_5}.
\tag{2.7.T8}\label{p1:eq:2.7.T8}
\]
The block sum is $D\sqrt R$ times a sum with weight
$\gamma^{-1}W_1S(m\bar r,\varsigma A_an;sc)$.
Put
\[
Z=\frac{4\pi\sqrt{A_aMN'}}{CS\sqrt R},
\qquad x=\frac{4\pi\sqrt{A_amn}}{cs\sqrt r}.
\tag{2.7.T9}\label{p1:eq:2.7.T9}
\]
Invert $c=4\pi\sqrt{A_amn}/(xs\sqrt r)$ and multiply by fixed smooth
cutoffs in $m,n,r,s$, each equal to 1 on that variable's range
$(y,2y]$ ($y=M,N',R,S$) and supported in $[y/2,3y]$. Since
$x/Z=\sqrt{(m/M)(n/N')(R/r)}\,(C/c)(S/s)$ and $C<c\le2C$, the
resulting smooth function $F(x,m,n,r,s)$ is supported on
$1/(12\sqrt3)\le x/Z\le6\sqrt2$, hence on $Z/21\le x\le9Z$.
Mellin inversion in its four log coordinates gives
\[
W_1=(2\pi)^{-4}\int_{\mathbb R^4}\phi_{\mathbf t}(x)
 m^{it_1}n^{it_2}r^{it_3}s^{it_4}\,d\mathbf t,
\tag{2.7.T10}\label{p1:eq:2.7.T10}
\]
\[
\max_{j\le5}\sup_x|x^j\phi_{\mathbf t}^{(j)}(x)|
 \ll\mathcal X^{17\epsilon_2}
                  \prod_{\nu=1}^4(1+|t_\nu|)^{-3}
 =:\mathfrak B(\mathbf t).
\tag{2.7.T11}\label{p1:eq:2.7.T11}
\]
For (2.7.T11), integrate three times by parts in each Mellin variable. There are four variables. Allow five $x$-derivatives. This uses at most 17 derivatives in (2.7.T8). The chain rule uses
$\partial_xc=-c/x,\partial_mc=c/(2m),\partial_nc=c/(2n)$,
$\partial_rc=-c/(2r),\partial_sc=-c/s$.
Thus (2.7.T10) and all ensuing integral interchanges are absolutely
convergent. For fixed $\mathbf t$ set
$\alpha_m=m^{it_1}1_{\mathcal I_j}(m)$ and
$\beta_{n,r,s}=b_{n,r,s}n^{it_2}r^{it_3}s^{it_4}$.
Estimate the following sum.
\[
\mathcal L_{\mathbf t}
 =\sum_{r,s,m,n}\alpha_m\beta_{n,r,s}
  \sum_{(c,r)=1}\gamma^{-1}\phi_{\mathbf t}
       (4\pi\sqrt{mA_an}/\gamma)S(m\bar r,\varsigma A_an;sc).
\tag{2.7.T12}\label{p1:eq:2.7.T12}
\]

\subsubsection{The cusp pair, both signs, and transform bounds}\label{p1:sec:t-cusp-pair}

Use $\sigma_\infty=I$ and
\[
\sigma_{1/s}=
 \begin{pmatrix}\sqrt r&0\\s\sqrt r&1/\sqrt r\end{pmatrix}.
\tag{2.7.T13}\label{p1:eq:2.7.T13}
\]
Multiply $\begin{pmatrix}a&b\\rsc_0&d\end{pmatrix}$ by (2.7.T13). With $c=rc_0+d$, the positive cusp moduli are $\gamma=cs\sqrt r$, $(c,r)=1$, as in (2.13).
Conversely choose $d\equiv c\ (r)$, $(d,sc)=1$, and
$c_0=(c-d)/r$; this supplies every allowed modulus. Double cosets
correspond to the reduced residues modulo $sc$. Insert their top and bottom entries in the cusp Kloosterman definition. This gives
\[
S_{\infty,1/s}(m,\nu;\gamma)
 =e(\nu\bar s/r)S(m\bar r,\nu;sc),
\qquad \nu\ne0.
\tag{2.7.T14}\label{p1:eq:2.7.T14}
\]
The numerator modulo $r$ is $\nu c$. Modulo $sc$ it is $m\bar r\,u+\nu\bar u$, with $u=r\bar d$.
This proves (2.7.T14) for either sign and for complex Fourier
coefficients.

We use the trace formula in DI's coefficient normalization, with the transforms $\widehat\phi,\check\phi,\widetilde\phi$ of (2.12).
Let $m,\nu>0$ and $r>R\ge1$. Use the same-sign Drappeau formula of Section 2. Apply the exact dictionary below.
\[
\rho^{\rm Dr}_{u,\mathfrak a}(n)=\frac{\rho^{\rm DI}_{u,\mathfrak a}(n)}
                                      {2\sqrt{|n|}},\quad
\rho^{\rm Dr}_{\mathfrak c\mathfrak a}(n,t)=
 \frac{\pi^{1/2+it}|n|^{it-1/2}}{\Gamma(1/2+it)}
                    \varphi_{\mathfrak c\mathfrak a n}(1/2+it).
\tag{2.7.T15}\label{p1:eq:2.7.T15}
\]
The holomorphic dictionary is
$\rho^{\rm Dr}_{f,\mathfrak a}(n)=(4\pi n)^{-\ell/2}
\psi_f(\mathfrak a,n)$.
For signs $+,-$, the Maass weights are as follows.
\[
h_+(\kappa)=-\widehat\phi(\kappa)/\cosh\pi\kappa,\quad
h_-(\kappa)=\frac\pi2\check\phi(\kappa)/\cosh\pi\kappa;
\tag{2.7.T16}\label{p1:eq:2.7.T16}
\]
For measure $dt$, the continuous weights are as follows.
\[
e_+(t)=-\widehat\phi(t)/\pi,\qquad e_-(t)=\check\phi(t)/2.
\tag{2.7.T17}\label{p1:eq:2.7.T17}
\]
Only sign $+$ has a holomorphic coefficient. It is
\[
\frac{i^\ell(\ell-1)!}{\pi(4\pi)^{\ell-1}}\,
                                 \widetilde\phi(\ell-1).
\tag{2.7.T18}\label{p1:eq:2.7.T18}
\]
There is no $1/\cosh(\pi t)$ in (2.7.T17).

For the opposite sign the normalization follows from the following derivation. Take a width-one scaling matrix $\sigma_{\mathfrak a}$. Let $\Gamma_{\mathfrak a}$ stabilize $\mathfrak a$ in $\Gamma=\Gamma_0(q)$, including $-I$. Then
$\sigma_{\mathfrak a}^{-1}\Gamma_{\mathfrak a}\sigma_{\mathfrak a}
=\{\pm T^k:k\in\mathbb Z\}$. Define, for $m>0$,
\[
 U_{\mathfrak a m}(z,s)=
 \sum_{\delta\in\Gamma_{\mathfrak a}\backslash\Gamma}
 (\Im\sigma_{\mathfrak a}^{-1}\delta z)^s
 e(m\sigma_{\mathfrak a}^{-1}\delta z),\qquad
 \langle F,G\rangle=\int_{\Gamma\backslash\mathbb H}
 F(z)\overline{G(z)}\,\frac{dx\,dy}{y^2}.
\]
Each coset enters once, with weight one. The series and
unfolding converge absolutely for $\Re s>1$.
Unfold the first series and decompose the second into cusp double
cosets of positive lower-left entry $\gamma$. Translation in the
real variable leaves the Kloosterman factor
$S_{\mathfrak a\mathfrak b}(m,-\nu;\gamma)$, with the
cusp convention of (2.7.T14). Explicitly, a matrix
$\sigma_{\mathfrak b}^{-1}\delta\sigma_{\mathfrak a}
=(\begin{smallmatrix}a&b\\\gamma&d\end{smallmatrix})$
gives phase $e((\nu a-md)/\gamma)$; its negative inverse
has entries $(-d,b;\gamma,-a)$, giving the same phase in
$S_{\mathfrak a\mathfrak b}(m,-\nu;\gamma)$.
The remaining integral is
\[
 I_\gamma=\gamma^{-2s_2}\int_{\mathbb H}
 y^{s_1+s_2-2}|z|^{-2s_2}
 \exp\!\left(2\pi i\{mz-\nu/(\gamma^2z)\}\right)dx\,dy.
\]
The zero-lower-left coset, if present for equivalent cusps,
vanishes by the integral of $e((m+\nu)x)$ over $[0,1]$.
Put $s_-=s_1-s_2$, $s_+=s_1+s_2$,
$z=i\varrho e^{i\alpha}$, $-\pi/2<\alpha<\pi/2$, and
$\varrho=(\sqrt{\nu/m}/\gamma)e^u$. Then
\[
 I_\gamma=\gamma^{-s_+}(\nu/m)^{s_-/2}
 \int_{-\pi/2}^{\pi/2}\cos^{s_+-2}\alpha
 \int_{\mathbb R}e^{s_-u}
 e^{-(4\pi\sqrt{m\nu}/\gamma)\cosh(u+i\alpha)}du\,d\alpha.
\]
For each interior $\alpha$, shift $u+i\alpha$ to the real
line. The vertical sides vanish exponentially, and the inner integral
is $2e^{-is_-\alpha}K_{s_-}(4\pi\sqrt{m\nu}/\gamma)$.
Let $\Re s_i>1$. At either angular endpoint the radial growth is of order $\cos^{- |\Re s_-|}\alpha$. For $\Re s_-=0$, a logarithm may also occur. The remaining exponent is $\Re s_+-2-|\Re s_-|=2\min\Re s_i-2>0$. Hence the original absolute integral is finite.
Hence Fubini and the endpoint limits are legitimate. Applying (2.7.T20) gives the following identity for $\Re s_1,\Re s_2>1$:
\[
\begin{aligned}
\langle U_{\mathfrak a m}(s_1),
             \overline{U_{\mathfrak b\nu}(s_2)}\rangle
={}&\pi2^{3-s_1-s_2}
 \frac{\Gamma(s_1+s_2-1)}{\Gamma(s_1)\Gamma(s_2)}
 (\nu/m)^{(s_1-s_2)/2}\\
&\times\sum_\gamma\gamma^{-s_1-s_2}
 S_{\mathfrak a\mathfrak b}(m,-\nu;\gamma)
 K_{s_1-s_2}(4\pi\sqrt{m\nu}/\gamma).
\end{aligned}
\tag{2.7.T19}\label{p1:eq:2.7.T19}
\]
The polar substitution and contour shift above give
the factor $2K_{s_1-s_2}$. The remaining angular integral is
\[
\int_{-\pi/2}^{\pi/2}\cos^{s_1+s_2-2}\alpha\,
 e^{-i(s_1-s_2)\alpha}\,d\alpha
 =\pi2^{2-s_1-s_2}
   \frac{\Gamma(s_1+s_2-1)}{\Gamma(s_1)\Gamma(s_2)}.
\tag{2.7.T20}\label{p1:eq:2.7.T20}
\]
For an independent normalization, unfold once and use (2.2). This gives
\[
\langle U_{\mathfrak a m}(s),u_j\rangle
 =\overline{\rho_{j\mathfrak a}(m)}
 \frac{\sqrt\pi(4\pi m)^{1/2-s}}{\Gamma(s)}
 \Gamma(s-\tfrac12-i\kappa_j)\Gamma(s-\tfrac12+i\kappa_j).
\]
This is the integral of $y^{s-3/2}e^{-2\pi my}K_{i\kappa_j}(2\pi my)$. Use the standard integral representation of $K$ and a beta integral. The second factor
$\langle u_j,\overline{U_{\mathfrak b\nu}(s_2)}\rangle$
is the same expression with $\rho_{j\mathfrak b}(-\nu)$,
$\nu,s_2$ in place of $\overline{\rho_{j\mathfrak a}(m)}$,
$m,s$. There is no conjugation of $s_2$ in this bilinear
pairing. Multiplication gives exactly (2.7.T21). Applying the same
unfolding to the Eisenstein expansion (2.4), with its standard
Parseval measure, gives (2.7.T22). Thus Parseval gives the spectral side with
\[
\begin{aligned}
P&=\frac{\pi m^{1/2-s_1}\nu^{1/2-s_2}(4\pi)^{1-s_1-s_2}}
          {\Gamma(s_1)\Gamma(s_2)},\\
\Lambda(s_1,s_2;t)&=
 \prod_{j=1}^2\Gamma(s_j-\tfrac12-it)\Gamma(s_j-\tfrac12+it).
\end{aligned}
\tag{2.7.T21}\label{p1:eq:2.7.T21}
\]
Its Maass summand is
$P\overline{\rho_{j\mathfrak a}(m)}
\rho_{j\mathfrak b}(-\nu)\Lambda(s_1,s_2;\kappa_j)$.
Its continuous integral is
\[
\begin{split}
&P\sum_{\mathfrak c}\int
 (\nu/m)^{it}\overline{\varphi_{\mathfrak c\mathfrak a m}(1/2+it)}
 \varphi_{\mathfrak c\mathfrak b,-\nu}(1/2+it)\\
&\qquad\times\Lambda(s_1,s_2;t)|\Gamma(1/2+it)|^{-2}\,dt.
\end{split}
\tag{2.7.T22}\label{p1:eq:2.7.T22}
\]
The factor at the second cusp has index $-\nu$. It arises because
the second Eisenstein parameter in Parseval is conjugated:
$\overline{E(z,1/2+it)}=E(z,1/2-it)$. So the continuous
product pairs a conjugated coefficient at $1/2+it$ with an unconjugated one at $1/2+it$.

For the unscaled Eisenstein series the continuous Parseval measure is
$dt/(4\pi)$. Its coefficient of
$\sqrt yK_{it}(2\pi|n|y)e(nx)$ is
\[
b_{(\mathfrak c,\mathfrak a)}(n,t)=
2\pi^{1/2+it}|n|^{it}\
\frac{\varphi_{\mathfrak c,\mathfrak a,n}(1/2+it)}
     {\Gamma(1/2+it)}.
\]
Hence $(4\pi)^{-1}\overline{b_{(\mathfrak c,\mathfrak a)}(m,t)}
b_{(\mathfrak c,\mathfrak b)}(-\nu,t)$ equals
\[
(\nu/m)^{it}
\overline{\varphi_{\mathfrak c,\mathfrak a,m}(1/2+it)}
\varphi_{\mathfrak c,\mathfrak b,-\nu}(1/2+it)
|\Gamma(1/2+it)|^{-2}.
\]
The coefficient product contributes $4\pi$, canceling the
Parseval denominator. Thus (2.7.T22) has plain $dt$ and no hyperbolic-cosine denominator. The second $\varphi$ is unconjugated at the same $t$.

Write $\sigma_i=\Re s_i>1$,
$\delta=|\sigma_1-\sigma_2|$,
$\sigma_*=\min(\sigma_1,\sigma_2)$, and
$x=4\pi\sqrt{m\nu}/\gamma$, with fixed $m,\nu\ge1$. Replace both Poincaré summands by their moduli. For one double coset with positive lower-left entry, the absolute integral is
\[
\ll_{\sigma_1,\sigma_2,m,\nu}
\gamma^{-\sigma_1-\sigma_2}
\bigl(1+\gamma^\delta+\log(2+\gamma)\bigr).
\tag{2.7.T23}\label{p1:eq:2.7.T23}
\]
The radial integral is a fixed positive factor times
$K_\delta(x\cos\alpha)$. For $\delta>0$, its small-argument
bound is $O_\delta((x\cos\alpha)^{-\delta})$; for $\delta=0$,
it is $O(1+|\log(x\cos\alpha)|)$. The angular weight is
$\cos^{\sigma_1+\sigma_2-2}\alpha$. Thus the power endpoint has
exponent $2\sigma_*-2>0$, while the logarithmic endpoint is
integrable and contributes $O(1+\log(2+\gamma))$ after angular
integration. The implied constant depends on the fixed spectral parameters and
on $m,\nu$, and is uniform in $\gamma$.

Consequently the controlling absolute-sum majorant is
\[
\sum_\gamma N_{\mathfrak a\mathfrak b}(\gamma)
\gamma^{-2\sigma_*}\bigl(1+\log(2+\gamma)\bigr)<\infty.
\tag{2.7.T24}\label{p1:eq:2.7.T24}
\]
For the cusp pair $(\infty,1/s)$, (2.7.T14) gives
$\gamma=cs\sqrt r$, $(c,r)=1$, and
$N(\gamma)=\varphi(sc)\le sc$. The sum is bounded by a constant
times
$\sum_{c\ge1}c^{1-2\sigma_*}(1+\log(2+c))$, which converges
because $\sigma_*>1$. For a general fixed cusp pair, the cited
cumulative count
$\sum_{\gamma\le T}N_{\mathfrak a\mathfrak b}(\gamma)
\ll_{\mathfrak a,\mathfrak b}T^2$ and dyadic decomposition bound
the tail by
$\sum_{k\ge0}(1+k)2^{-2(\sigma_*-1)k}<\infty$. Any bounded
initial range is finite.

The logarithm in (2.7.T23) is attained: for $\sigma_1=\sigma_2=2$, $m=\nu=1$ and the admissible moduli $\gamma=\sqrt2 c$, $c$ odd, the absolute cell integral is $\pi\gamma^{-4}\log\gamma+O(\gamma^{-4})$. The summable majorant (2.7.T24) supplies the uniform step beyond the fixed-$\gamma$ angular endpoint estimate.

Both sides continue to the connected domain
$\Re s_i>3/4,\ |\Re(s_1-s_2)|<1/2$. For the geometric side,
Weil's bound and $K_\mu(x)\ll x^{-|\Re\mu|}(1+|\log x|)$
give locally uniform absolute convergence. For the spectral side, apply Stirling and DI's positive large sieves. These give locally uniform absolute convergence. Since $i\kappa_j\le7/64$, the exceptional gamma factors have no poles here. Equality therefore
holds at $s_1=1+iu,s_2=1-iu$. With
\[
\mathcal H(t,u)=\frac{\cosh\pi t}
                 {\cosh\pi(t-u)\cosh\pi(t+u)},
\tag{2.7.T25}\label{p1:eq:2.7.T25}
\]
this gives
\[
\begin{aligned}
&\sum_\gamma\frac{4\pi\sqrt{m\nu}}{\gamma^2}
 S(m,-\nu;\gamma)K_{2iu}(4\pi\sqrt{m\nu}/\gamma)\\
&\qquad=\frac{\pi^2}2\sum_j\frac{\bar\rho_j(m)\rho_j(-\nu)\mathcal H(\kappa_j,u)}
                      {\cosh\pi\kappa_j}
 +\frac\pi2\sum_{\mathfrak c}\int E_{\mathfrak c}(t)\mathcal H(t,u)\,dt,
\end{aligned}
\tag{2.7.T26}\label{p1:eq:2.7.T26}
\]
where
\[
 E_{\mathfrak c}(t)
 :=(\nu/m)^{it}
   \overline{\varphi_{\mathfrak c\mathfrak a m}(1/2+it)}
   \varphi_{\mathfrak c\mathfrak b,-\nu}(1/2+it).
\]
The factor $|\Gamma(1/2+it)|^{-2}$ of (2.7.T22) is part of the
spectral kernel; $E_{\mathfrak c}$ excludes it.
The constant eigenfunction contributes $0$ to the sum over $j$: its
Fourier coefficients vanish at the nonzero frequencies $m$ and $-\nu$.

To integrate (2.7.T26) against a test function, set
\[
\mathcal R(u)=\int_0^\infty K_{2iu}(x)\phi(x)x^{-2}\,dx,\qquad
k(u)=\frac4{\pi^2}u\sinh(2\pi u)\mathcal R(u).
\tag{2.7.T27}\label{p1:eq:2.7.T27}
\]
Kontorovich--Lebedev inversion gives
$\int_\mathbb R k(u)K_{2iu}(x)\,du=\phi(x)/x$;
see \cite[\S10.43(v), (10.43.30)--(10.43.31)]{DLMF}.
Use the following kernel identity.
\[
\int_\mathbb R u\sinh(2\pi u)\mathcal H(t,u)K_{2iu}(x)\,du
       =\cosh(\pi t)xK_{2it}(x).
\tag{2.7.T28}\label{p1:eq:2.7.T28}
\]
For $|\Re\mu|+|\Re\nu|<1$, the absolutely convergent product integral is
\[
\int_0^\infty K_\mu(x)K_\nu(x)\,dx
 =\frac{\pi^2}
 {4\cos(\pi(\mu+\nu)/2)\cos(\pi(\mu-\nu)/2)}.
\tag{2.7.T29}\label{p1:eq:2.7.T29}
\]
Use the Bessel product representation
$K_\mu K_\nu(x)=2\int_0^\infty K_{\mu+\nu}(2x\cosh v)
\cosh((\mu-\nu)v)\,dv$. Integrate in $x$. For $|\Re b|<1$, use
$\int_0^\infty\cosh(bv)/\cosh v\,dv=(\pi/2)\sec(\pi b/2)$.
The latter follows by expanding $1/\cosh v$ and the partial
fractions of secant. In the subsequent KL application,
$(\mu,\nu)=(2it,2iu)$ with $u\in\mathbb R$; here $t\in\mathbb R$ on
the regular and continuous spectrum, while an exceptional parameter
satisfies $|\Im t|\le 7/64$. Thus $|\Re\mu|+|\Re\nu|\le 7/32<1$ and
$|\Re(\mu-\nu)|<1$. Apply KL inversion to
$xK_{2it}(x)$, whose KL transform is, by (2.7.T29),
$(\pi^2/4)\mathcal H(t,u)/\cosh\pi t$. This proves (2.7.T28).

All interchanges here are justified as follows. Split the integral representation of $K_{iv}$ at $\sinh U=\max(1,1/x)$. Integrate by parts on its tail. This gives
\[
|K_{iv}(x)|\ll e^{-\pi|v|/2}
                    (1+|v|+\log(2+1/x)).
\tag{2.7.T30}\label{p1:eq:2.7.T30}
\]
Apply the self-adjoint operator $(x\partial_x)^2-x^2$ repeatedly to $\phi(x)/x$. Then
$\mathcal R(u)\ll_{\phi,j}e^{-\pi|u|}(1+|u|)^{1-2j}$.
On real $t$, $|\mathcal H(t,u)|\le
4\cosh(\pi t)e^{-2\pi\max\{|t|,|u|\}}$; exceptional parameters
give an exponentially decaying function of $u$ as well.
DI's positive quadratic estimates imply convergence after the
remaining $(1+|t|)^{-4}$ bound. The geometric side is absolutely
summable by Weil. These bounds justify Fubini on both sides of
(2.7.T26). Equations (2.7.T27)--(2.7.T28) now give exactly
(2.7.T16)--(2.7.T17) for sign $-$. For sign $+$, let $r>R\ge1$. Insert (2.7.T15) in the classical same-sign formula. This gives the minus signs in (2.7.T16)--(2.7.T17) and the coefficient (2.7.T18). These are the
normalizations of Section 2, (2.14).

Let $\phi$ be supported in $[Z/21,9Z]$. With $\mathfrak B=\max_{j\le5}\sup x^j|\phi^{(j)}(x)|$, the transform bounds are
\[
|\widehat\phi(t)|+|\check\phi(t)|
 \ll\mathfrak B(1+|\log Z|)(1+Z)^{-1}
       \min\{1,((1+Z)/|t|)^3\}\quad(t\in\mathbb R),
\tag{2.7.T31}\label{p1:eq:2.7.T31}
\]
\[
|\widehat\phi(i\nu)|+|\check\phi(i\nu)|
 \ll\mathfrak B(1+|\log Z|)(1+Z^{-{2\nu}})/(1+Z),
\quad0<\nu\le1/4,
\tag{2.7.T32}\label{p1:eq:2.7.T32}
\]
and the same bound as (2.7.T31) for $\widetilde\phi(\ell)$ with
$|t|$ replaced by $\ell$. To prove them, put
$\Phi(y)=\int\cos(xy)\phi(x)\,dx/x$ and
$\Phi_c(y)=\int e^{-ixy}\phi(x)\,dx/x$.
Five integrations by parts give
$|\Phi^{(j)}(y)|+|\Phi_c^{(j)}(y)|
\ll\mathfrak B Z^j\min(1,(Z|y|)^{-5})$, $j\le3$.
The exact Bessel representations are
\[
\begin{aligned}
\check\phi(t)&=\frac4\pi\int_0^\infty
  \cos(2tu)\Phi(\sinh u)\,du,\\
\widehat\phi(t)&=-2\int_0^\infty
  \cos(2tu)\Phi(\cosh u)\,du,\\
\widetilde\phi(\ell)&=\frac1{2\pi}\int_{-\pi}^{\pi}
  e^{i\ell u}\Phi_c(\sin u)\,du.
\end{aligned}
\]
Their absolute bounds follow by splitting at $\sinh u=1/Z$;
three integrations by parts give the cubic decay. The periodic
holomorphic integral has no boundary terms. Replacing cosine by
$\cosh(2\nu u)$ and splitting at the same point gives (2.7.T32).
Truncate the oscillatory representations first. For real $t$, the
tails are $O((1+|t|)/(x\sinh U))$. For $t=i\nu$, they are
$O((1+\nu)e^{2\nu U}/(x\sinh U))$. Since $2\nu\le1/2$, these
tails tend to zero and justify the interchanges.

\subsubsection{Cusp transport and all Hecke factors}\label{p1:sec:t-cusp-transport}

Choose integers $w,y$ with $rw-sy=1$, and let
\[
W=r^{-1/2}\begin{pmatrix}r&y\\rs&rw\end{pmatrix}.
\tag{2.7.T33}\label{p1:eq:2.7.T33}
\]
Then
$\sigma_{1/s}^{-1}W=\begin{pmatrix}1&y/r\\0&1\end{pmatrix}$,
with $y/r\equiv-\bar s/r(1)$. Also $W\Gamma_0(rs)W^{-1}
=\Gamma_0(rs)$. Here is an integer computation of the latter fact.
For
$A=r^{-1/2}\begin{pmatrix}r\alpha&\beta\\rs\gamma&r\delta\end{pmatrix}$,
$r\alpha\delta-s\beta\gamma=1$, conjugating
$\begin{pmatrix}a&b\\rsc&d\end{pmatrix}$ gives entries
\[
\begin{pmatrix}
 a\alpha\delta r-\alpha b\gamma rs+\beta c\delta rs-\beta d\gamma s&
 -a\alpha\beta+\alpha^2br+\alpha\beta d-\beta^2cs\\
 rs(a\delta\gamma-b\gamma^2s+c\delta^2r-d\delta\gamma)&
 -a\beta\gamma s+\alpha b\gamma rs+\alpha d\delta r-\beta c\delta rs
\end{pmatrix}.
\tag{2.7.T34}\label{p1:eq:2.7.T34}
\]
Both $W=A(1,y,1,w)$ and $W^{-1}=A(w,-y,-1,1)$ have this
form, proving both inclusions. Thus composition by $W$ is unitary,
commutes with the Laplacian, and preserves each eigenspace. The
weight-$\ell$ slash action has the same properties.

Comparison of Fourier series after the translation $y/r$ gives
\[
\rho_{F,1/s}(\nu)=e(\nu\bar s/r)c_{F\circ W}(\nu),\quad
\psi_F(1/s,\nu)=e(\nu\bar s/r)\psi_{F|_\ell W}(\infty,\nu).
\tag{2.7.T35}\label{p1:eq:2.7.T35}
\]
For Eisenstein series, substitute $W\gamma W^{-1}$ in the
defining coset series for $\Re z>1$. It gives
$E_{\mathfrak c}(Wz,u)=E_{W^{-1}\mathfrak c}(z,u)$, using
$W^{-1}\sigma_{\mathfrak c}$ as a width-one scaling matrix.
Changing the source-cusp representative or translating its scaling
matrix does not change the Fourier coefficient at the fixed
observation cusp $\infty$.
Meromorphic continuation therefore gives
\[
\varphi_{\mathfrak c,1/s,\nu}(u)
 =e(\nu\bar s/r)\varphi_{W^{-1}\mathfrak c,\infty,\nu}(u).
\tag{2.7.T36}\label{p1:eq:2.7.T36}
\]
In (2.7.T12), (2.7.T14) contributes
$e(-\varsigma A_an\bar s/r)$. Equations (2.7.T35)--(2.7.T36)
cancel this phase exactly for both signs, in every spectrum.

For $(k,q)=1$, define the normalized Hecke operator below.
\[
T_kF(z)=k^{-1/2}\sum_{ad=k}\sum_{b\bmod d}F((az+b)/d).
\tag{2.7.T37}\label{p1:eq:2.7.T37}
\]
Sum each Fourier term, of index $\nu'$, over $b$. This forces
$d\mid\nu'$, and the term moves to index $\nu=a\nu'/d$. The
coefficient at $\nu\ne0$ is
\[
c_{T_kF}(\nu)=\sum_{a\mid(k,\nu)}c_F(k\nu/a^2).
\tag{2.7.T38}\label{p1:eq:2.7.T38}
\]
The prefactor is 1 because, in the normalization (2.2) of
\cite[(1.15)]{DI82}, the factor $\sqrt y$ contributes $\sqrt{a/d}$,
and $k^{-1/2}d\sqrt{a/d}=1$.
For weight $\ell$, use
$T_k=k^{(\ell-1)/2}\sum_{ad=k}d^{-\ell}\sum_b(\cdot)$;
the normalized coefficients
$a_F(\nu)=\nu^{-(\ell-1)/2}\psi_F(\infty,\nu)$
satisfy the same identity at coprime $(k,\nu)$.

Split $n=n'n''$, $n''\mid a_T^\infty$, $(n',a_T)=1$,
and put $k=A_an''$. Then $(k,q)=(k,n')=1$.
The number of pieces is
\[
\mathcal D\le(2N_1)^\epsilon
                  \prod_{p\mid a_T}(1-p^{-\epsilon})^{-1}
 \ll_\epsilon\mathcal X^{2\epsilon},
\quad \tau(k)^2\ll_\epsilon\mathcal X^\epsilon.
\tag{2.7.T39}\label{p1:eq:2.7.T39}
\]
The product bound is uniform in $a_T$. Choose $p_\epsilon$ so that $(1-p^{-\epsilon})^{-1}\leq p^\epsilon$ for $p>p_\epsilon$. The smaller primes contribute a constant depending only on $\epsilon$. The remaining product is at most $|a_T|^\epsilon$.
On a Maass eigenspace $V$, the transported $n''$-functional is
$Y\circ T_k\circ W$, where
$Y(v)=\sum_{n'}\beta_{n'n''}c_v(\varsigma n')$.
By Riesz representation and unitarity,
\[
\sum_{u_j\in V}|\ell_{2,n''}(u_j)|^2
 \le\tau(k)^2k^{\theta_{\mathrm{KS}}}\sum_{u_j\in V}|Y(u_j)|^2.
\tag{2.7.T40}\label{p1:eq:2.7.T40}
\]
The self-adjoint unramified $T_k$ diagonalizes on the full
level space. By Kim--Sarnak its norm is at most
$\tau(k)k^{\theta_{\mathrm{KS}}/2}$. For holomorphic forms replace $k^{\theta_{\mathrm{KS}}}$
by 1, by Deligne. For negative index, $c_v(-n')=\overline{c_{\bar v}(n')}$. Conjugation permutes an orthonormal basis. Hence the positive-index quadratic bound holds with conjugated coefficients. The argument allows complex
coefficients and any orthonormal basis of each eigenspace.

The continuous spectrum requires its own proof. Represent a cusp by
$u/w$, $w\mid q$, put $w'=(w,q/w)$. The orbit conditions in
\cite[Lemma~3.6, (3.9)--(3.11)]{DI82}
$(c_1,d_1)=1,(c_1,q)=w,c_1d_1\equiv uw\pmod{ww'}$, with $c_1=w\gamma$,
give exactly
\[
\begin{aligned}
\varphi_{u/w,\infty,n}(z)&=
 (w'/(wq))^z
 \sum_{\gamma\ge1,(\gamma,q/w)=1}\gamma^{-2z}\\
&\qquad
 \sum_{\substack{d_1\bmod\gamma w\\(d_1,\gamma w)=1,\ \gamma d_1\equiv u(w')}}
                         e(-nd_1/(\gamma w)).
\end{aligned}
\tag{2.7.T41}\label{p1:eq:2.7.T41}
\]
For a character $\chi\bmod w'$, let
$\mathcal V_\chi(n;z)=\sum_{u\bmod w'}^*\chi(u)\varphi_{u/w,\infty,n}(z)$.
For $(k,q)=(k,n')=1$, either sign of $n'$, the identity is
\[
\mathcal V_\chi(kn';z)=\bar\chi(k)
          \left(\sum_{d\mid k}\chi^2(d)d^{1-2z}\right)\mathcal V_\chi(n';z).
\tag{2.7.T42}\label{p1:eq:2.7.T42}
\]
To prove it, sum (2.7.T41) over $u$. Write $c_1=c_1'c_1''$, where
$c_1''$ is supported on primes of $q$ and $(c_1',q)=1$.
CRT factors the inner sum as
$\chi(c_1')G_\chi(-n;c_1'')c_{c_1'}(n)$.
The outer character supplies another $\chi(c_1')$.
Replacing the Gauss-sum variable by $\bar kd_1$ gives
$\bar\chi(k)$ on $c_1''$. The other factor is
\[
L^{(q)}(2z,\chi^2)^{-1}
     \sum_{d\mid|n|,(d,q)=1}\chi^2(d)d^{1-2z},
\]
by $c_{c_1'}(n)=\sum_{d\mid(c_1',n)}d\mu(c_1'/d)$; this divisor sum
factors at $kn'$. All manipulations converge absolutely for
$\Re z>1$. Both sides are meromorphic Fourier coefficients,
so (2.7.T42) holds meromorphically, hence at $z=1/2+it$ outside
a discrete set. There are finitely many cusp characters.
Orthogonality in $u\bmod w'$ gives
\[
\begin{split}
&\sum_{\mathfrak c}\left|\sum_{n'}v_{n'}|kn'|^{it}
            \varphi_{\mathfrak c,\infty,kn'}(1/2+it)\right|^2\\
&\qquad\le\tau(k)^2\sum_{\mathfrak c}\left|\sum_{n'}v_{n'}|n'|^{it}
            \varphi_{\mathfrak c,\infty,n'}(1/2+it)\right|^2.
\end{split}
\tag{2.7.T43}\label{p1:eq:2.7.T43}
\]
Here the character eigenvalue is
$\bar\chi(k)\sum_{d\mid k}\chi^2(d)(k/d^2)^{it}$,
of modulus at most $\tau(k)$. Transport (2.7.T36) permutes the
complete cusp set. For negative indices,
$\varphi_{\mathfrak c\infty,-n'}(1/2+it)
=\overline{\varphi_{\mathfrak c\infty,n'}(1/2-it)}$.
Changing $t$ to $-t$ and conjugating coefficients converts
each symmetric integration window to DI's positive-index
continuous quadratic form. The discrete exceptional set is null
for the $dt$-integral.

\subsubsection{Regular, continuous, and holomorphic estimates}

Put $Q_T=RS$, $L_{n''}=N'/n''\ge1/2$. Include $|t|<1$ among the spectral windows $K$. By (2.7.T31), the undivided transform satisfies the following bound.
\[
\mathfrak h(K)\ll\mathfrak B(\mathbf t)\mathcal X^{o(1)}
 (1+Z)^{-1}\min\{1,((1+Z)/K)^3\}.
\tag{2.7.T44}\label{p1:eq:2.7.T44}
\]
For Maass forms use the positive DI weights $1/\cosh\pi\kappa_j$. For holomorphic forms use the positive weights $(\ell-1)!/(4\pi)^{\ell-1}$. These weights are split equally
between the two Cauchy--Schwarz factors. For the continuous spectrum use plain $dt$. Take the supremum of $e_\varsigma(t)$ outside the positive quadratic form.
By (2.7.T17), $|e_\varsigma|\le|\Phi_\varsigma|/2$, with
$\Phi_+=\widehat\phi,\Phi_-=\check\phi$. Thus every norm uses a
positive measure.

\cite[Theorem~2]{DI82} at $\infty$ gives the $m$-side
\[
\sum_{q\asymp Q_T}\tau(q)(K^2+M^{1+\epsilon}/q)M
 \ll\mathcal X^{3\epsilon}(Q_TK^2+M)M.
\tag{2.7.T45}\label{p1:eq:2.7.T45}
\]
For the $n$-side, Cauchy--Schwarz over $n''$ costs $\mathcal D$. Apply (2.7.T40) or (2.7.T43). Then use \cite[Theorem~2]{DI82} at scale $L_{n''}$. The bound is
\[
\ll\mathcal D\sum_{n''}\tau(k)^2k^{\theta_{\mathrm{KS}}}
             (K^2+L_{n''}^{1+\epsilon}/Q_T)\|b^\flat(n'')\|_2^2.
\tag{2.7.T46}\label{p1:eq:2.7.T46}
\]
On continuous and holomorphic parts (2.7.T46) holds with $k^{\theta_{\mathrm{KS}}}$
replaced by 1. \cite[Theorem~2]{DI82} allows the endpoint $M=1/2$
in (2.7.T45).

Use $K^2+U\le((1+Z)^2+U)\max(1,K/(1+Z))^2$.
The sum over dyadic $K$ of the latter factor times the cubic
decay in (2.7.T44) is $O(1+\log(2+Z))$. The cubic decay also absorbs a factor $K^\epsilon$: with it the sum is $O((1+Z)^\epsilon\log(2+Z))$. Thus the bound (2.6) in the form $(K^2+N/q)(KN)^\epsilon$ gives the same estimate. Cauchy--Schwarz and
(2.7.T39) therefore yield
\[
\begin{aligned}
|\mathcal L_{\mathbf t}^{\rm reg}|
\ll{}&\mathcal X^{4\epsilon}\mathfrak B(\mathbf t)(1+Z)^{-1}
 [(Q_T(1+Z)^2+M)M]^{1/2}\\
&\times\left[\sum_{n''}(A_an'')^{\theta_{\mathrm{KS}}}
 ((1+Z)^2+L_{n''}/Q_T)\|b^\flat(n'')\|_2^2\right]^{1/2}.
\end{aligned}
\tag{2.7.T47}\label{p1:eq:2.7.T47}
\]
The same calculation with plain $dt$ proves (2.7.T47) for
the continuous part. It uses the symmetry shown after
(2.7.T43). The same cubic summation proves absolute convergence
of its integral, justifying the window partition. Thus the
continuous part uses the measure $dt$ throughout.

\subsubsection{Exceptional spectrum and the one-integer block}\label{p1:sec:t-exceptional}

Write $L=L_{n''}$, $\Xi=1/Z$, $\Omega=1+Q_T/L$.
Since $0<\vartheta_j\le\theta_{\mathrm{KS}}$,
\[
1+\Xi^{\vartheta_j}
 \le2(1+\Xi/\sqrt{\Omega})^{\vartheta_j}\Omega^{\vartheta_j/2}.
\tag{2.7.T48}\label{p1:eq:2.7.T48}
\]
Fix $n''$. Apply Cauchy--Schwarz, assigning
$Y=(1+\Xi/\sqrt{\Omega})^2$ to the $m$-side and
$\Omega$ to the $n$-side. With $\ell_{2,n''}$ as in (2.7.T40), the two exceptional quadratic forms are
\[
\begin{aligned}
\mathcal M_{n''}&=\sum_{r,s}\sum_{\rm exc}^{(rs)}Y^{\vartheta_j}
 \Bigl|\sum_m\alpha_m\rho_{j\infty}(m)\Bigr|^2,\\
\mathcal N_{n''}&=\sum_{r,s}\sum_{\rm exc}^{(rs)}\Omega^{\vartheta_j}
 |\ell_{2,n''}(u_j)|^2,
\end{aligned}
\]
with $r,s$ over the support of $b$. Let $L\ge1/2$. Apply (2.7.T40) and Theorem 2.2 at level $rs>Q_T$, length $L$ and $Y=\Omega$. Since $L\Omega/(rs)\le1+L/Q_T$, the factor in (2.8) is $\ll\max(1,L)^\epsilon(1+L/Q_T)$. Summing over $r,s$ gives
\[
\mathcal N_{n''}\ll\mathcal X^{2\epsilon}(A_an'')^{\theta_{\mathrm{KS}}}
                 (1+L/Q_T)\|b^\flat(n'')\|_2^2.
\tag{2.7.T49}\label{p1:eq:2.7.T49}
\]
Put $Q'=4Q_T$, $Y_0(M)=(Q'+M)^2/M$, and
\[
Z_{n''}=\frac{C^2S^2R}{(Q_T+M)^2(L+Q_T)}.
\tag{2.7.T50}\label{p1:eq:2.7.T50}
\]
The identity
$\Xi^2/\Omega=(16\pi^2A_an'')^{-1}
C^2S^2R/[M(L+Q_T)]$ gives
\[
\max(1,Y/Y_0(M))^{\theta_{\mathrm{KS}}}
 \le1+(A_an'')^{-\theta_{\mathrm{KS}}}Z_{n''}^{\theta_{\mathrm{KS}}}.
\tag{2.7.T51}\label{p1:eq:2.7.T51}
\]
Let $M\ge1$. Then (2.9) applies to every integer interval
inside $(M,2M]$. Use the positive Hilbert space with weights $Y^{\vartheta_j}$. Apply Abel summation to $m^{it_1}$. The norm is at most $1+|t_1|\log2$ times the maximum interval norm. Take the maximum after the norm. It remains outside the spectral sum. Apply (2.9) at $Y_0(M)$. Then use monotonicity
\[
Y^{\vartheta_j}\le\max(1,Y/Y_0)^{\theta_{\mathrm{KS}}} Y_0^{\vartheta_j},
\tag{2.7.T52}\label{p1:eq:2.7.T52}
\]
For fixed $q$, the coefficient $\alpha_m$ is independent of the
factorization $q=rs$. Put
$w(q)=\#\{(r,s)\text{ in the current support}:rs=q\}$ and
$\Delta=\max_{2\le q\le Q'}w(q)$, and let $\omega_0(q)$
count the distinct prime factors of $q$. Since $(r,s)=1$,
$w(q)\le2^{\omega_0(q)}\le\tau(q)\ll_\eta q^\eta$.
In the Hilbert space over all levels $2\le q\le Q'$ and their exceptional
forms, vector-valued Abel summation gives
$\mathcal M_{n''}^{1/2}\le\Delta^{1/2}
(1+|t_1|\log2)\sup_{I\subset(M,2M]}
F_Y(I)^{1/2}$, where
$F_Y(I)=\sum_{2\le q\le Q'}\sum_{\rm exc}^{(q)}Y^{\vartheta_j}
|\sum_{m\in I}\rho_{j\infty}(m)|^2$ and $I$ ranges over integer
intervals. Thus the supremum is outside the whole quadratic form.
Apply (2.9) with exponent $\eta=\epsilon/2$ at $Y_0(M)$.
By (2.7.T4), $M\le H\le4\mathcal X^2$ when $\epsilon_1\le1$;
also $Q'=4Q_T\le4\mathcal X$. Consequently
$\Delta(Q'M)^\eta\ll_\epsilon\mathcal X^{2\epsilon}$.
Together with (2.7.T51)--(2.7.T52), this is absorbed in
$\mathcal X^{3\epsilon}$ and gives
\[
\mathcal M_{n''}\ll\mathcal X^{3\epsilon}(1+t_1^2)
 [1+(A_an'')^{-\theta_{\mathrm{KS}}}Z_{n''}^{\theta_{\mathrm{KS}}}](Q_T+M)M.
\tag{2.7.T53}\label{p1:eq:2.7.T53}
\]

\cite[(8.17)]{DI82} excludes $M=1/2$. We prove that endpoint here.
Let $p(q)$ be the least prime not dividing $q$. Every smaller
prime divides $q$. The classical Chebyshev lower bound
$\Theta(x):=\sum_{\ell\le x}\log\ell\gg x$ for $x\ge2$
(see \cite{IK04}) gives
\[
p(q)\ll1+\log q.
\tag{2.7.T54}\label{p1:eq:2.7.T54}
\]
Here $\Theta(p(q)-1)\le\log q$, because the product of all
smaller primes divides $q$; the case $p(q)=2$ is immediate.
The implied constant in (2.7.T54) is absolute and independent of $a_T$.

On an exceptional eigenspace $V$, let $v$ represent
$u\mapsto c_u(1)$. Self-adjointness and (2.7.T38) show that
$T_pv,T_{p^2}v$ represent $c_u(p),c_u(p^2)$.
The Hecke identity $T_p^2=T_{p^2}+I$, with
$\|T_p\|\le2p^{\theta_{\mathrm{KS}}/2}$, gives
\[
\sum_{u_j\in V}|c_{u_j}(1)|^2
 \le8p^{\theta_{\mathrm{KS}}}\sum_{u_j\in V}|c_{u_j}(p)|^2
       +2\sum_{u_j\in V}|c_{u_j}(p^2)|^2.
\tag{2.7.T55}\label{p1:eq:2.7.T55}
\]
Multiply by $Y^{\vartheta_j}$, which is constant on $V$. Group
levels by $p(q)\le P:=C_0(1+\log(4Q_T))$, where $C_0$ is an
absolute constant supplied by (2.7.T54). Positivity permits
enlarging each group to all $2\le q\le Q'$. For $k=p,p^2$, the singleton $\{k\}$ is an integer interval in $(k/2,k]$, with $k/2\ge1$, so (2.9) applies at $N=k/2$. In DI's terms, subtract the initial intervals $(k/2,k]$ and $(k/2,k-1]$. The resulting factor 4 is included in $c_\epsilon$.

For the endpoint, the same factorization count multiplies the
level-$q$ quadratic form by at most $\tau(q)$. The following bound
holds for any positive exponent. In its use in (2.7.T57), take
$\eta=\min(\epsilon/2,1/4)$ and use
$\tau(q)\ll_\epsilon Q_T^{\epsilon/2}$ for
$Q_T<q\le4Q_T$. This keeps their combined loss within the
displayed $\epsilon$-power after adjustment of its constant. At
$Y_k=(Q'+k/2)^2/(k/2)$ it gives
\[
\sum_{2\le q\le Q'}\sum_{\rm exc}^{(q)}
 Y^{\vartheta_j}|c_{u_j}(k)|^2
 \ll_\epsilon (Q'k)^\epsilon(Q'+k/2)k
                          \max(1,Y/Y_k)^{\theta_{\mathrm{KS}}}.
\tag{2.7.T56}\label{p1:eq:2.7.T56}
\]
Since $Y_0(1/2)/Y_k=k(Q'+1/2)^2/(Q'+k/2)^2\le k\le9k/4$, insert (2.7.T56), with its
exponent taken as $\eta$, in (2.7.T55).
Sum over the at most $P$ possible primes:
\[
\begin{aligned}
&\sum_{Q_T<q\le4Q_T}\tau(q)\sum_{\rm exc}^{(q)}
 Y^{\vartheta_j}|c_{u_j}(1)|^2\\
&\qquad\ll_\epsilon Q_T^\epsilon P^6
 (Q_T+\tfrac12)\tfrac12\,
 \max(1,Y/Y_0(1/2))^{\theta_{\mathrm{KS}}},
\end{aligned}
\tag{2.7.T57}\label{p1:eq:2.7.T57}
\]
for $0<\epsilon\le1/4$, with the implied constant including
the divisor bound and the change from $\epsilon$ to $\eta$.
Before simplification, the sum is at most a constant times
\[
c_\epsilon(9/4)^{\theta_{\mathrm{KS}}} Q'^{1+\epsilon}
 \sum_{p\le P}(8p^{2+2\theta_{\mathrm{KS}}+\epsilon}
                         +2p^{4+2\theta_{\mathrm{KS}}+2\epsilon});
\]
this is $\ll_\epsilon Q_T^{1+\epsilon}P^6$, because
$5+2\theta_{\mathrm{KS}}+2\epsilon\le6$.
The loss $P^6$ is logarithmic and uniform in $a_T$.
In particular $P^6\ll_\epsilon\mathcal X^\epsilon$; together
with the factorization loss already included in (2.7.T57), this
fits the $\mathcal X^{3\epsilon}$ allowance of (2.7.T53).
Together with (2.7.T51), this proves (2.7.T53) also for $M=1/2$.
Since $\alpha_1=1$, Abel summation costs nothing here.

Apply the triangle inequality over $n''$ before Cauchy--Schwarz.
Then (2.7.T49), (2.7.T53), and (2.7.T32) give
\[
\begin{aligned}
|\mathcal L_{\mathbf t}^{\rm exc}|
\ll{}&\mathcal X^{6\epsilon}\mathfrak B(\mathbf t)(1+|t_1|)
 (1+Z)^{-1}\sqrt{M(Q_T+M)}\\
&\times\sum_{n''}\big[(A_an'')^{\theta_{\mathrm{KS}}/2}
                 +Z_{n''}^{\theta_{\mathrm{KS}}/2}\big]
        ((L+Q_T)/Q_T)^{1/2}\|b^\flat(n'')\|_2.
\end{aligned}
\tag{2.7.T58}\label{p1:eq:2.7.T58}
\]
There is no exceptional continuous or holomorphic spectrum.

\subsubsection{The four terms and completion of the proof}

Put
\[
\mathfrak T=A_aMN'/(C^2S),\quad M_0=CS/D,\quad
\varpi=4\mathcal X^{\epsilon_2}.
\tag{2.7.T59}\label{p1:eq:2.7.T59}
\]
The quantity $\mathfrak T$ is unrelated to the cell endpoint $T$ of
Section~3. Then $M\le\varpi M_0$,
$Z^2Q_T=16\pi^2\mathfrak T$, so
$(1+Z)^{-2}Q_T^{-1}\asymp(Q_T+\mathfrak T)^{-1}$.
Square (2.7.T47) and multiply by $D^2R$. Apart from the stated auxiliary powers, sum the following over $n''$.
\[
(A_an'')^{\theta_{\mathrm{KS}}} D^2R
 \frac{M(M+Q_T+\mathfrak T)(L+Q_T+\mathfrak T)}
                       {Q_T+\mathfrak T}\|b^\flat(n'')\|_2^2.
\tag{2.7.T60}\label{p1:eq:2.7.T60}
\]
The first, noncancelling branch of (2.7.T58) has the same bound. Termwise, use
\[
(Q_T+M)(L+Q_T)/Q_T
 \le[Q_T(1+Z)^2+M][(1+Z)^2+L/Q_T].
\tag{2.7.T61}\label{p1:eq:2.7.T61}
\]
To obtain (2.42), use the following identity.
\[
\frac{Q_T(M+Q_T+\mathfrak T)(L+Q_T+\mathfrak T)}
                            {Q_T+\mathfrak T}
=(L+Q_T)(M+Q_T)+Q_T\mathfrak T
                         -\frac{LM\mathfrak T}{Q_T+\mathfrak T}.
\tag{2.7.T62}\label{p1:eq:2.7.T62}
\]
Multiply (2.7.T62) by $M/Q_T$, then drop the last (nonpositive) term:
\[
\frac{M(M+Q_T+\mathfrak T)(L+Q_T+\mathfrak T)}{Q_T+\mathfrak T}
\le\frac{M(L+Q_T)(M+Q_T)}{Q_T}+M\mathfrak T.
\]
Since
$D^2M/S\le\varpi CD$ and
$M+Q_T\le\varpi S(C+DR)/D$, it follows that
\[
\begin{aligned}
&D^2R\frac{M(M+Q_T+\mathfrak T)(L+Q_T+\mathfrak T)}
                            {Q_T+\mathfrak T}\\
&\qquad\le\varpi^2
 \{CS(L+RS)(C+DR)+A_aN'RS\}.
\end{aligned}
\tag{2.7.T63}\label{p1:eq:2.7.T63}
\]
The last term follows from
$D^2RM\mathfrak T=A_aN'RS(DM/(CS))^2$.
Using $L\le N_1/n''$, $N'\le N_1$, $A_a\le A$ proves
the first two terms of (2.42), with
$A^{\theta_{\mathrm{KS}}}(n'')^{\theta_{\mathrm{KS}}}$ kept explicit.

For the cancelling branch of (2.7.T58), use
\[
Z_{n''}^{\theta_{\mathrm{KS}}}(L+Q_T)
 \le(C^2S^2R)^{\theta_{\mathrm{KS}}}(Q_T+M)^{-2\theta_{\mathrm{KS}}}(N_1+Q_T)^{1-\theta_{\mathrm{KS}}},
\tag{2.7.T64}\label{p1:eq:2.7.T64}
\]
and $\sum_{n''}\|b^\flat(n'')\|_2
\le\mathcal D^{1/2}\|b\|_2$. Since
$M\mapsto M(Q_T+M)^{1-2\theta_{\mathrm{KS}}}$ is increasing,
substitute $M_0$, paying at most $\varpi^{2-2\theta_{\mathrm{KS}}}$.
The resulting identity is
\[
\begin{aligned}
&\frac{D^2RM_0(Q_T+M_0)^{1-2\theta_{\mathrm{KS}}}
                  (C^2S^2R)^{\theta_{\mathrm{KS}}}(N_1+Q_T)^{1-\theta_{\mathrm{KS}}}}{Q_T}\\
&\qquad=CS(CD\sqrt R)^{2\theta_{\mathrm{KS}}}
             (N_1+RS)^{1-\theta_{\mathrm{KS}}}(C+DR)^{1-2\theta_{\mathrm{KS}}}.
\end{aligned}
\tag{2.7.T65}\label{p1:eq:2.7.T65}
\]
This is the third term of (2.42). The fourth was (2.7.T3).
The truncated tail is smaller than it.

To account for all losses, reconstruct the block sum as
$D\sqrt R(2\pi)^{-4}\int\mathcal L_{\mathbf t}\,d\mathbf t$.
Equation (2.7.T11) gives
$\int\mathfrak B(\mathbf t)(1+|t_1|)\,d\mathbf t
\ll\mathcal X^{17\epsilon_2}$.
The amplitude's $\varpi$ loss in (2.7.T63) is at most
$\mathcal X^{\epsilon_2}$ after taking its square root.
The log partitions and divisor factors cost $\mathcal X^{O(\epsilon_1)}$. This includes the DI small $\epsilon$ losses, $\mathcal D$, and the one-integer $P^6$ loss. The bound is uniform in $a_T$.
Consequently the final loss is
\[
\mathcal X^{O(\epsilon_1)+18\epsilon_0}.
\tag{2.7.T66}\label{p1:eq:2.7.T66}
\]
Choose the freely available $\epsilon_1$ small enough relative
to the requested $\epsilon$. Equations (2.7.T63), (2.7.T65), and (2.7.T3)
then prove (2.41)--(2.43), with loss
$\mathcal X^{\epsilon+18\epsilon_0}$, for both signs and
including the $M=1/2$ block. This proves Theorem 2.8 in the form used in Section~3, (3.14)--(3.15).

